\documentclass[11pt,a4paper]{article}
\usepackage[T1]{fontenc}
\usepackage[utf8]{inputenc}
\usepackage{lmodern,microtype}
\usepackage[margin=26mm]{geometry}
\usepackage{amsmath,amssymb,amsthm,mathtools}
\numberwithin{equation}{section}
\usepackage{booktabs,longtable,tabularx,array}
\newcolumntype{P}[1]{>{\raggedright\arraybackslash}p{#1}}
\usepackage{enumitem,xcolor,fancyhdr}
\usepackage{listings}
\usepackage[hidelinks,breaklinks]{hyperref}
\usepackage{xurl}
\hypersetup{pdftitle={From finite observations to continuous structures},pdfauthor={Research report prepared for Vladimir Reshetnikov}}
\definecolor{ink}{HTML}{20344A}
\definecolor{accent}{HTML}{176B75}
\definecolor{pale}{HTML}{F0F5F7}
\pagestyle{fancy}\fancyhf{}
\fancyhead[L]{\small\color{ink}Continuous structures over iPRA}
\fancyhead[R]{\small 20 November 2026}
\fancyfoot[C]{\thepage}
\setlength{\headheight}{14pt}
\setlength{\parindent}{0pt}
\setlength{\parskip}{5pt plus 1pt}
\setlist{itemsep=3pt,topsep=4pt}
\emergencystretch=2em
\newtheorem{theorem}{Theorem}[section]
\newtheorem{lemma}[theorem]{Lemma}
\newtheorem{proposition}[theorem]{Proposition}
\newtheorem{corollary}[theorem]{Corollary}
\theoremstyle{definition}
\newtheorem{definition}[theorem]{Definition}
\newtheorem{example}[theorem]{Example}
\theoremstyle{remark}\newtheorem{remark}[theorem]{Remark}
\newcommand{\N}{\mathbb N}\newcommand{\Z}{\mathbb Z}
\newcommand{\Q}{\mathbb Q}\newcommand{\R}{\mathbb R}\newcommand{\C}{\mathbb C}
\newcommand{\eps}{\varepsilon}\newcommand{\ipra}{\ensuremath{\mathsf{iPRA}}}
\newcommand{\PR}{\ensuremath{\mathsf{PR}}}\newcommand{\dist}{\operatorname{d}}
\newcommand{\norm}[1]{\left\lVert #1\right\rVert}
\newcommand{\abs}[1]{\left|#1\right|}
\newcommand{\clamp}{\operatorname{clip}}
\newcommand{\e}[1]{2^{-(#1)}}
\newcommand{\code}[1]{\texttt{#1}}
\newcommand{\sourcebase}{https://github.com/VladimirReshetnikov/ProveIt/blob/5f3be02388442f16487d2dc5c2529961dd783d2b/Logic/PrimitiveRecursiveArithmetic/}
\lstset{basicstyle=\ttfamily\footnotesize,breaklines=true,columns=fullflexible,keepspaces=true,frame=single,rulecolor=\color{pale}}
\begin{document}
\hypersetup{pageanchor=false}
\begin{titlepage}
\color{ink}
{\Large Research report\par}
\vspace{18mm}
{\huge\bfseries From finite observations\\[4pt]to continuous structures\par}
\vspace{7mm}
{\Large Quantitative completions, compact function spaces,\\continuous model semantics, and local-to-global analysis\\over intuitionistic primitive recursive arithmetic\par}
\vspace{13mm}
{\large Prepared for Vladimir Reshetnikov\par}
{\large 20 November 2026\par}
\vspace{12mm}
\begin{minipage}{0.94\textwidth}
\textbf{Main result.} Explicit presentations by finite rational observations support continuous objects, maps between their completions, compact spaces of functions, and continuous logical quantifiers. The arithmetic layer proves the finite inequalities; an explicit adequacy argument explains their continuous meaning. Coherence, witnesses, domain margins, and primitive-recursive bounds remain visible throughout.
\end{minipage}
\vfill
\small
\textbf{Source baseline.} ProveIt commit
\path{5f3be02388442f16487d2dc5c2529961dd783d2b}.
\par\textbf{Verification boundary.} The paper gives complete human arguments and a detailed object-theory formalization path. A standalone Lean kernel checks finite extrema and continuous-formula transfer. The new continuous theorems have not all been reified as \code{Prf PRAAx} derivations. Exact checked declarations and remaining obligations are listed in the verification chapter.
\end{titlepage}
\hypersetup{pageanchor=true}
\begin{abstract}
The obstruction to using continuous mathematics in \ipra\ is rarely just the absence of real-number objects in its language. The harder issues are obtaining coherent approximations, controlling quantification over infinite domains, and transporting local certificates without silently adding choice, comprehension, or a universal evaluator. We give a layered solution for a substantial but explicitly delimited fragment.

First, uniformly continuous finite evaluators extend to completed metric presentations, and a quantitative uniqueness modulus turns separately computed approximate solutions into a canonical exact solution. Second, positive overlap weights glue local approximations with an explicit global continuity bound. Third, finite clipped McShane envelopes give primitive-recursive nets consisting of actual bounded Lipschitz functions. This makes function spaces available as compact presented sorts. Fourth, a total approximate correspondence transfers continuous formulas with an error independent of quantifier depth for nonexpansive connectives. Thus nested suprema and infima can be evaluated by finite tables without constructing a truth predicate. We give an end-to-end variational example and a singular inverse example.

These are conservative numerical realizations for fixed presentations and fixed formulas, not a claimed interpretation of unrestricted analysis or model theory. We identify which components already exist in ProveIt and which interfaces this report completes. Complete Lean verification of the finite formula-transfer kernel accompanies detailed definitions, arithmetic invariants, and a plan for reification into the repository's ordinary iPRA calculus.
\end{abstract}
\begingroup
\small
\setlength{\parskip}{0pt}
\tableofcontents
\endgroup
\clearpage
\section{The target theory and the actual gap}\label{sec:foundation}

\subsection{Scope: continuous mathematics with a numerical interface}

Our target is first-order intuitionistic primitive recursive arithmetic: equality, a symbol for each fixed primitive-recursive definition, its defining equations, the successor axioms, and induction for quantifier-free formulas. Formulas may have quantifiers; the restriction is on induction. The ordinary target in ProveIt is a derivation
\[
  []\ \vdash[\mathsf{PRAAx}]\ \varphi,
\]
where \(\varphi\) is a concrete first-order formula over its primitive-recursive terms. A Lean theorem about natural numbers, rational numbers, or a metric space is not by itself a derivation of this form.

The objective is to use continuous mathematical objects without pretending that \ipra\ has a sort containing every real number, every map, every category, or every model. A compact metric space can be accessed by finite nets; a real-valued function by rational samples and a modulus; a function space by finite nets of genuine functions; and a continuous formula by a rational value with a certified error. All the finite computations are arithmetic. The assertion that they faithfully describe a familiar continuous object is a separate, proved adequacy statement.

This distinction is productive. We will quantify over spaces of functions in the continuous semantics, while the arithmetic statement quantifies only over natural numbers encoding a finite net and a table of observations. We will obtain exact continuous solutions from error bounds at every precision, while never using an algorithm that tests equality of arbitrary real numbers.

\subsection{What the current source already supplies}

The source audit is pinned to \cite{repo}; it is not a claim to have rebuilt the repository. Its directory overview substantially understates the newer analytic source tree. The following interfaces were inspected directly.

\begin{longtable}{@{}P{.28\textwidth}P{.32\textwidth}P{.33\textwidth}@{}}
\toprule
Source under \path{Derivations/Lean/IPRADerivations/} & Existing result & Boundary addressed here\\
\midrule\endhead
\path{Foundations/CoherentMinimizer.lean} & Quantitative uniqueness gives a Cauchy name for selected approximate minimizers. & Identify the value of the completed point and prove completed uniqueness through a general evaluation-extension theorem.\\
\path{Foundations/StrongMinimizer.lean} & Finite-net argmin plus a strong-convexity certificate yields a distance bound. & Use an arbitrary residual-diameter modulus, including a singular inverse with a H\"older modulus.\\
\path{Foundations/FiniteGluing.lean} & Exact compatible finite tables can be glued by choosing a patch. & Glue noisy local continuous evaluators using weights and an explicit positive denominator bound.\\
\path{Foundations/FiniteNet.lean} & Arithmetic estimates transfer a finite minimum along a supplied covering/modulus interface. & Construct a net of valid elements of an infinite-dimensional function space.\\
\path{Foundations/RealizerMasterInterface.lean} & An actual finite list of input certificates is packaged as a derived interface, with soundness and stability derivations. & Supply analytic and model-theoretic providers for such interfaces rather than propose another generic wrapper.\\
\path{Foundations/CoherenceObstruction.lean} & An object-level diagonal step for a finite-branch obstruction. Its general computability conclusion is separately identified. & State exactly which positive coherence datum defeats the obstruction in our applications.\\
\path{Foundations/UniversalEvaluatorObstruction.lean} & A fixed-program diagonal theorem and a separate external no-universal-PR-evaluator theorem. & Compile each fixed presentation; do not evaluate all PR programs through one PR function.\\
\path{Foundations/FiniteUseSignObstruction.lean} & An external Lean theorem shows that finite observations cannot decide exact real sign. & Use continuous truth values and robust inequalities, retaining the exact-sign boundary.\\
\bottomrule
\end{longtable}

The prior \path{docs/proofs/foundations/ipra_embedding_framework} packet already contains quantitative implicit-function arguments and finite mesh positivity. The current \path{Analysis/GridIntegral.lean} and \path{Analysis/RiemannRefinement.lean} also supply arithmetic integration estimates. Consequently, neither ``use a modulus'' nor ``replace a continuum by a mesh'' is advertised as a new idea. The contribution is the combination of a completed-value interface, a constructive function-space presentation, and continuous-formula transfer, with their exact strength and remaining proof obligations stated.

\subsection{Three levels, with different quantifiers}

\begin{definition}[Three meanings of a theorem]\label{def:levels}
An \emph{arithmetic theorem} is an \ipra\ sentence involving fixed PR programs and numerical codes. A \emph{presentation schema} is a metatheoretic construction that, for every fixed tuple of PR programs and supplied \ipra\ proof values, produces such a sentence and derivation. A \emph{semantic adequacy theorem} says what the numerical theorem means for completed metric spaces or other continuous objects.
\end{definition}

The schema quantifier ``for every fixed program'' is external. It is not an object quantifier over a universal evaluator. An oracle-relative algorithm can describe a mathematical construction before the input evaluator is fixed. It becomes an unconditional PR program only after substituting fixed PR input programs, or after replacing all oracle calls by lookups in a supplied finite table. Calling the oracle-relative algorithm primitive recursive must always retain that qualification.

For example, if \(a(k)\) and \(b(k)\) are fixed PR rational sequences, the regularity assertion
\[
 \forall k\,\forall l\quad |a(k)-a(l)|\leq 2^{-k}+2^{-l}
\]
is an ordinary first-order arithmetic sentence after clearing denominators. Quantifying instead over every arbitrary sequence \(a:\N\to\Q\) belongs to the schema or the semantics. We never identify those two readings.

\subsection{Codes and bounded predicates}

Here is a concrete coding convention sufficient to remove all hidden data types. Use
\(\langle a,b\rangle=2^a(2b+1)-1\) for pairs. Its two inverse coordinates are obtained by searches bounded by the input plus one, so they are PR. A list is a length together with a nested pair code; list access is a bounded iteration and returns a designated default out of range. Tuples are fixed-length lists. Integers are differences of naturals. A rational is a signed numerator and a positive natural denominator; distinct codes may denote the same rational. Equality and order are decided by cross multiplication, without requiring a canonical fraction.

All well-formedness tests, finite sums and products, rational minima and maxima, dyadic powers, rounding of rational numbers, and scans of a list are PR. A failed positive-denominator test returns an error flag and a default rational. Every claimed output theorem retains the well-formedness guard; default values never turn malformed data into a certificate.

A bounded numerical statement is represented by its characteristic function. For example,
\[
  \bigwedge_{i<b}P(i)\quad\hbox{means}\quad
  \mathsf{all}_P(b)=1,
\]
where \(P\) is decidable by a fixed PR bit. Bounded existential statements are treated by a finite search or a stored witness. Their correspondence with the displayed bounded formulas is proved by quantifier-free induction on the relevant scan. Consequently, an invariant containing finite bounded quantifiers is legitimate only after this translation. Merely describing an invariant in ordinary mathematical prose does not grant an unrestricted induction principle.

\subsection{What a presented continuous object contains}

\begin{definition}[Metric presentation]\label{def:presentation}
A metric presentation has numerical codes for a dense family of basic points, a decidable well-formedness predicate, and a rational distance evaluator \(\Delta(a,b,p)\). Its finite contracts prove the nonnegativity, symmetry, and triangle inequalities with the error associated with precision \(p\), and prove that the distance samples form a regular name. Equivalently, in its semantic interpretation it has a metric \(d\) such that
\[
 |\Delta(a,b,p)-d(a,b)|\leq 2^{-p}.
\]
The equivalence is an adequacy theorem; the finite contracts, not this semantic equation alone, are the arithmetic input. An exact rational metric is the special case in which \(\Delta\) is independent of \(p\).

A point name gives basic-point codes at each precision with a proved Cauchy bound. A map presentation gives finite rational/basic-point output approximations, an input modulus, a range guard, and stability under equivalent input names. A compact presentation additionally gives a nonempty PR finite net at every scale and a covering contract. Completeness and total boundedness are separate data.
\end{definition}

For distance approximations, a useful finite list of laws is
\begin{align*}
 |\Delta(a,b,p)-\Delta(a,b,q)|&\leq 2^{-p}+2^{-q},\\
 \Delta(a,c,p)&\leq\Delta(a,b,p)+\Delta(b,c,p)+3\cdot2^{-p},\\
 |\Delta(a,b,p)-\Delta(b,a,p)|&\leq2\cdot2^{-p}.
\end{align*}
Add nonnegativity up to \(2^{-p}\), a vanishing self-distance bound, and a specified zero-distance equivalence. Sending the error to zero gives the corresponding pseudometric on basic points; quotienting by zero distance is only semantic notation. The arithmetic development uses the setoid relation and its finite approximations. Function spaces below naturally use this approximate-distance form, even though their basic functions are evaluated exactly on rational inputs.

\begin{definition}[Observation contract]\label{def:observation}
A finite query consists of a code \(q\), a precision \(k\), and numerical parameters \(u\). An input table \(T\) records the answers to all required subqueries. A PR predicate \(\mathsf{Good}(T,q,k,u)\) checks its finite shape, domain guards, demanded compatibility inequalities, and all retained certificates. A fixed PR transformer \(A\) produces the output. Its arithmetic obligation is
\begin{equation}
 \forall T,q,k,u\quad
 \mathsf{Good}(T,q,k,u)\ \longrightarrow\
 \mathsf{Out}(A(T,q,k,u),q,k,u).
 \label{eq:contract}
\end{equation}
An independent provider theorem constructs good tables from a chosen presentation. An independent adequacy theorem interprets \(\mathsf{Out}\) as the desired continuous estimate.
\end{definition}

Checking a finite table is not evidence that every later table exists or that successive choices are coherent. Formula \eqref{eq:contract} deliberately cannot conceal either issue. The completion and residual theorems below supply coherence in specific cases. The compact-function-space theorem supplies actual net centers, avoiding a hidden search for a function satisfying a quantized sample table.

\subsection{A precise finite-realization principle}

\begin{theorem}[Finite arithmetic realization]\label{thm:realization}
Fix a finite collection of PR numerical primitives, with \ipra\ derivations of their required finite contracts. Form a program using composition, finite tuples, rational arithmetic with checked denominators, decidable branches, finite extrema, and primitive recursion. For every recursive loop supply a quantifier-free encoded invariant, derivations of its initial and step clauses, and a derivation from its final clause to the output contract. Then the program has a PR numerical implementation and an \ipra\ derivation of \eqref{eq:contract}. Replacing each closed primitive by its definition gives a theorem in the original language of \ipra.
\end{theorem}

\begin{proof}
Composition substitutes PR terms and cuts the constituent derivations. Tuples and rational operations are defined by the fixed coding routines. A branch condition is a PR equality or comparison; its decidability is proved inside arithmetic, and intuitionistic disjunction elimination combines the branch derivations. A recursive loop is a PR iterator on its coded state. Apply quantifier-free induction to its encoded invariant, then substitute the actual loop bound. A finite nesting of these constructions uses only a finite nesting of PR recursion. Finally expand the definitions of the finitely many derived symbols. Each step is an intuitionistic derivation transformation, so no appeal to the semantic truth of arbitrary \ipra\ proofs occurs.

For a minimum scan, store both the running minimum \(m_j\) and its index \(r_j\). The invariant says \(r_j\leq j\), \(m_j=a_{r_j}\), and \(m_j\leq a_i\) for every \(i\leq j\). The last assertion is its bounded characteristic bit. On inserting \(a_{j+1}\), use the decidable comparison to retain \((m_j,r_j)\) or replace it by \((a_{j+1},j+1)\). Both branch implications are finite ordered arithmetic. This gives attained minima and all the extrema estimates used later without induction on an unbounded quantified formula. Maxima are analogous.
\end{proof}

The theorem is a deliberately explicit proof-construction principle, not a theorem that arbitrary terminating algorithms or arbitrary true inequalities are provable in \ipra. The substantive work in each application is to prove its local estimates and give the invariants. Our later sections do that at the human level; the Lean plan identifies where each is reified. When a local estimate is only an assumption, it remains an assumption in the resulting arithmetic theorem.

\subsection{Errors, equality, and the strength of the interpretation}

Use \(\eps_k=2^{-k}\). If one execution returns an approximation within \(\eps_k\) and another within \(\eps_l\) of the same semantic value, their outputs differ by at most \(\eps_k+\eps_l\). Conversely, a proved bound of that form constructs a regular real name without first assuming a semantic value. Many ``exact'' conclusions below mean equality of such names: every requested rational error is met.

This does not give a Boolean equality test. A robust inequality \(x<r\) can be witnessed by a sample whose upper error endpoint is below \(r\). Equality is an all-precision property and generally has no finite deciding sample. Division requires an explicit positive lower bound on the denominator, and branch choices based on sign require an explicit gap or a continuous replacement.

The conservative claim is correspondingly limited. New arithmetic notation expands into PR terms and derived rules, so the numerical theory is a definitional extension of \ipra. We do not assert that unrestricted real analysis, ordinary model-theoretic compactness, all categories, or a higher-order system with arbitrary choice is conservative over \ipra. Even classical conservation results over PRA must be read in their stated theories and formula classes; they do not automatically transfer to an intuitionistic language with witnesses \cite{kohlenbach}.

\providecommand{\anE}{\varepsilon}
\providecommand{\anD}{\mathcal D}
\providecommand{\anH}{\widehat}

\section{Continuous completion as a reusable interface}
\label{sec:continuous-completion}

The present section isolates three operations that turn finite analytic
estimates into continuous mathematics: extending a map to a completion,
gluing local maps with explicit weights, and identifying a completed point
from an approximate solution procedure.  The distinction between these
operations matters.  A Cauchy estimate by itself does not identify the
completed value of a function, and a collection of correct local answers
does not by itself give a continuous global answer.

There is already substantial relevant work in the inspected repository.
\path{CoherentMinimizer.cauchy} derives a regular name from a quantitative
uniqueness modulus; its source explicitly leaves canonical evaluation and
uniqueness of the completed point outside that theorem.
\path{StrongMinimizer.selected_pair_stable} obtains a distance estimate
for actual finite-net minimizers from a strong-convexity certificate.
\path{FiniteGluing.glue_restrict} glues exactly agreeing finite tables by
the first applicable patch.  The earlier analytic embedding packet also
constructs simple analytic branches under a derivative separation
certificate.  We reuse these principles.  The additions here are the
completion interface that finishes the value-identification step, a
quantitative continuous gluing theorem with approximate data, and a
singular inverse example that needs a H\"older modulus rather than a
nonzero derivative.

\subsection{Metric presentations and the scope of the statements}

Put $\anE_n=2^{-n}$.  For the exact rational core, a metric presentation
consists of a decidable set $D$ of natural-number codes and a fixed
primitive-recursive rational function $d:D\times D\to\Q_{\ge0}$ satisfying
symmetry, $d(x,x)=0$, and the triangle inequality.  Distinct codes may have
distance zero.  All membership and metric laws have explicitly specified
arithmetic formulas; an application must prove them or retain them as
premises.  We do not derive a metric law from the fact that its checker
terminates.

A regular name is a sequence $x=(x_n)_{n\ge0}$ in $D$ with
\begin{equation}
 d(x_m,x_n)\le\anE_m+\anE_n.
 \label{eq:an-regular}
\end{equation}
For internal use the sequence is a fixed primitive-recursive evaluator
$a(p,n)$, where $p$ is an ordinary numerical parameter.  Quantifying over
all names is external notation.  Each arithmetic theorem is a schema in
the fixed evaluator and its supplied contracts; it does not evaluate an
arbitrary code for a primitive-recursive program.

The semantic statements below also hold for arbitrary metric spaces.
For a general effective presentation, the distance itself has rational
approximants $\Delta(x,y,t)$ with error at most $\anE_t$.  Then
$d(x,y)\le q$ abbreviates
\[
 \forall t\quad \Delta(x,y,t)\le q+\anE_t.
\]
The exact rational version makes the arithmetic kernel easiest to see.
The extension to distance names is addressed after the main construction.
No induction is performed on a statement quantifying over real points.

\begin{lemma}[Removing a vanishing rational error]
\label{lem:an-small-error}
For rational numbers $u,v$, if $u\le v+C\anE_n$ for every natural $n$,
where $C\ge0$ is rational, then $u\le v$.  The contradiction index can be
chosen primitive recursively from the rational codes of $u,v,C$.
\end{lemma}

\begin{proof}
Rational order is decidable.  In the case $u>v$, write $u-v=a/b$ with
positive integers $a,b$, so $u-v\ge1/b$.  Choose an integer $K\ge C$ and
put $n=Kb+1$.  The elementary inequality $2^n> Kb$ gives
$C\anE_n<1/b$ when $K>0$, contradicting the assumed inequality.  If
$K=0$, the contradiction is immediate.  Integer arithmetic, bounded
rounding, and the proof of $2^n\ge n+1$ use quantifier-free induction.
The universal premise is instantiated once, at the displayed numerical
index; it is not the induction formula.
\end{proof}

Define the represented distance bound
\begin{equation}
 \anH d(x,y)\le q
 \quad\Longleftrightarrow\quad
 \forall m,n\;
 d(x_m,y_n)\le q+\anE_m+\anE_n,
 \qquad q\in\Q_{\ge0}.
 \label{eq:an-completion-order}
\end{equation}
Write $x\sim y$ when this holds with $q=0$.  Reflexivity and symmetry
follow immediately.  For transitivity, insert an intermediate sample
$y_t$ and use the two hypotheses to obtain
$d(x_m,z_n)\le\anE_m+\anE_n+2\anE_t$; then apply
Lemma~\ref{lem:an-small-error}.  The same proof gives the triangle
inequality for the represented distance bounds.  Thus regular names
modulo $\sim$ form the represented completion $\anH D$.  Here
completeness is asserted for Cauchy families with a supplied joint family
of names, as in Theorem~\ref{thm:an-diagonal}.  No choice of representatives
for an arbitrary abstract sequence of quotient points is inferred.
Its distance is
represented by the rational samples
\[
 d(x_{n+2},y_{n+2}).
\]
The reverse triangle inequality bounds their pairwise difference by
$2\anE_{m+2}+2\anE_{n+2}\le\anE_m+\anE_n$.
Constant names embed $D$ into $\anH D$, and
\begin{equation}
 \anH d(x,\operatorname{const}(x_n))\le\anE_n.
 \label{eq:an-density}
\end{equation}
Indeed, the defining inequalities are precisely regularity, with an
additional nonnegative error allowance for the constant name.

\begin{theorem}[Explicit diagonal completion]
\label{thm:an-diagonal}
Let $x_{i,p}\in D$ be a fixed two-parameter family such that every
row $p\mapsto x_{i,p}$ is a regular name, and suppose
\[
 d(x_{i,p},x_{j,q})
 \le\anE_i+\anE_j+\anE_p+\anE_q
 \quad\text{for all }i,j,p,q.
\]
The diagonal $z_n=x_{n+2,n+2}$ is a regular name, and its completed
distance from the $j$th name is at most $\anE_j$.  In particular, a
regular sequence of regular names with this explicit joint evaluator
can be flattened primitive recursively.
\end{theorem}

\begin{proof}
Substitute $(i,p)=(m+2,m+2)$ and $(j,q)=(n+2,n+2)$ to obtain
$d(z_m,z_n)\le2\anE_{m+2}+2\anE_{n+2}$, which is at most the
regularity budget.  For the distance to the $j$th name, substitute only
$(i,p)=(m+2,m+2)$:
$d(z_m,x_{j,q})\le2\anE_{m+2}+\anE_j+\anE_q
\le\anE_m+\anE_j+\anE_q$.
This is exactly~\eqref{eq:an-completion-order}.  There is one fixed
joint evaluator in the hypothesis; the proof never chooses a name for
each member of an arbitrary sequence of completed points.
\end{proof}

\subsection{Extension, continuity, and uniqueness}

Let $f:D\to E$ be a fixed primitive-recursive map between rational
metric presentations.  A dyadic uniform-continuity modulus is a function
$\mu:\N\to\N$ satisfying
\begin{equation}
 d_D(u,v)\le\anE_{\mu(k)}
 \quad\Longrightarrow\quad
 d_E(f(u),f(v))\le\anE_k.
 \label{eq:an-map-modulus}
\end{equation}
Replace $\mu(k)$ by $\max_{j\le k}\mu(j)$ to make it nondecreasing.
This replacement is primitive recursive, and its maximum invariant is
quantifier free.  Define
\begin{equation}
 h(n)=\mu(n+2)+2,
 \qquad (\mathsf{Ext}_f x)_n=f(x_{h(n)}).
 \label{eq:an-extension}
\end{equation}

\begin{theorem}[Quantitative universal extension]
\label{thm:an-extension}
Formula~\eqref{eq:an-extension} produces a regular name, respects
$\sim$, and extends $f$ on constant names.  The induced map
$\anH f:\anH D\to\anH E$ has modulus
\[
 \nu(k)=\mu(k+2)+2.
\]
It is the unique uniformly continuous extension of $f$ to these
completions.  All sample constructions and index changes are primitive
recursive when the input routines and moduli are fixed primitive
recursive functions.
\end{theorem}

\begin{proof}
If $m\ge n$, monotonicity of $h$ and regularity give
\[
 d_D(x_{h(m)},x_{h(n)})
 \le2\anE_{h(n)}
 =\tfrac12\anE_{\mu(n+2)}.
\]
The modulus therefore bounds the output distance by $\anE_{n+2}$,
which is stronger than the required regularity inequality.

Suppose now $\anH d_D(x,y)\le\anE_{\nu(k)}$.  For $m,n\ge k$,
the corresponding input distance is at most
\[
 \anE_{\nu(k)}+\anE_{h(m)}+\anE_{h(n)}
 \le 3\anE_{\mu(k+2)+2}
 \le\anE_{\mu(k+2)}.
\]
Consequently the distance between these output samples is at most
$\anE_{k+2}$.  Insert such sufficiently late samples between any two
fixed output samples, use their regularity, and let their indices tend
to infinity by Lemma~\ref{lem:an-small-error}.  This proves that the
completed output distance is at most $\anE_{k+2}$, hence at most
$\anE_k$.  If $x\sim y$, the same argument applies for every $k$ and
proves output equivalence.  Constant names are preserved by inspection.

For uniqueness, let $g$ be another extension with a supplied uniform
modulus $\lambda$.  Fix a name $x$ and an output precision $k$.
Choose
$n\ge\max(\nu(k+2),\lambda(k+2))$.  By
\eqref{eq:an-density}, both $\anH f(x)$ and $g(x)$ are within
$\anE_{k+2}$ of the common value $f(x_n)$.  Their mutual distance is
at most $2\anE_{k+2}$ for every $k$, so they are equivalent.
This argument uses a maximum of two known indices, not a search for a
continuity modulus.
\end{proof}

Identity maps and composition are preserved by this extension, as
equalities of represented maps.  For composition, both candidate maps
are uniformly continuous and agree on every constant name; uniqueness
applies.  Equality of the raw output programs is neither needed nor
generally true.  This is the useful categorical universal property of
completion at the interface of the arithmetic and semantic layers.

\paragraph{Distance names.}
The proof used only triangle inequalities, uniform-continuity
implications, and explicit positive slack.  With named distances,
replace every final bound $d\le q$ by all rational tests
$\Delta_t\le q+\anE_t$.  When an intermediate inequality contains a
sum of $r$ distance observations, evaluate each to a precision $t$ for
which $r\anE_t$ is below the allocated slack.  The displayed factors
$1/2$ and $3/4$ provide such slack; non-strict limiting bounds are
closed with Lemma~\ref{lem:an-small-error}.  The query index is computed
by bounded doubling from the number of terms and requested precision.
For real-valued semantic distances, this lemma is applied to the
rational distance tests, not directly to an unencoded real quantity.
The triangle and modulus laws for $\Delta$ are still explicit
hypotheses.  This gives the same extension schema for effective metric
presentations, including function spaces whose sup distances have only
certified rational approximants.

\section{Quantitative gluing of continuous local data}
\label{sec:analytic-gluing}

Exact membership in a chart is often a poor interface for a represented
point near the chart boundary.  For a convex target, a finite weighted
construction removes the need to select one chart.  Unlike exact
finite-table gluing, it also remains quantitatively stable when the
local values are only approximately compatible.

\subsection{Data and the rational algorithm}

Let $D$ be a rational metric presentation, let $c_1,\ldots,c_N\in D$,
and let $r_i>0$ be rational.  Define
\[
 w_i(x)=\max(0,r_i-d(x,c_i)),
 \qquad S(x)=\sum_{i=1}^N w_i(x).
\]
Fix a presented subdomain $D_0\subseteq D$ and assume a rational
$\delta>0$ satisfies $S(x)\ge\delta$ for every valid code in $D_0$.
All local contracts below are restricted to $D_0$, and every regular
input name used in this section has all its samples in $D_0$.
Thus the completed domain is $\widehat{D_0}$.  Membership of the limit
alone would not justify denominator bounds for the sampled codes.
The implementation retains the ambient metric presentation and a
separate predicate for $D_0$, so centers may lie outside $D_0$ while
remaining valid ambient points.  For each $i$ there is a fixed rational-vector
evaluator $f_i(p,x)\in\Q^d$.  The finite family is either fixed
externally or given by one fixed indexed evaluator; indices do not
encode arbitrary programs.  It is required to be correct only when
$w_i(x)>0$.  Use the $\ell^1$ norm
$\norm{v}_1=\sum_{j<d}|v_j|$, which is an exact rational finite sum.
Assume the finite local contracts
\begin{align}
 \norm{f_i(p,x)}_1&\le M,
 && w_i(x)>0,
 \label{eq:an-local-bound}\\
 \norm{f_i(p,x)-f_i(q,y)}_1
 &\le Ld(x,y)+\anE_p+\anE_q,
 && w_i(x),w_i(y)>0,
 \label{eq:an-local-lip}
\end{align}
where $M,L\ge0$ are rational.  These formulas include precision
coherence by taking $x=y$.  The targets can be real vectors, or complex
numbers represented by pairs of rationals; no complex order is used.

Set a summand to zero when its weight is zero, and otherwise evaluate
the local routine.  Define
\begin{equation}
 A_p(x)=\sum_i w_i(x)f_i(p,x),
 \qquad U_p(x)=\frac{A_p(x)}{S(x)}.
 \label{eq:an-glue}
\end{equation}
This is a total rational algorithm: division has a harmless zero
default outside the certified domain, and the positive denominator
contract proves that the default is never used on that domain.  No
local routine is required to satisfy an analytic assertion outside its
guard.  With exact rational distances the weight comparison is simply
a decidable rational comparison.

\begin{theorem}[Weighted gluing with sample errors]
\label{thm:an-glue}
The algorithm satisfies
\begin{align}
 \norm{U_p(x)}_1&\le M,\label{eq:an-glue-bound}\\
 \norm{U_p(x)-U_q(y)}_1
 &\le C d(x,y)+\anE_p+\anE_q,
 \qquad C=L+\frac{2NM}{\delta}.
 \label{eq:an-glue-lip}
\end{align}
If $j$ is active at $x$ and every active $i$ satisfies
\begin{equation}
 \norm{f_i(p,x)-f_j(q,x)}_1
 \le\eta+\anE_p+\anE_q,
 \label{eq:an-overlap}
\end{equation}
then
\begin{equation}
 \norm{U_p(x)-f_j(q,x)}_1
 \le\eta+\anE_p+\anE_q.
 \label{eq:an-glue-restrict}
\end{equation}
The error does not acquire a factor $N$ in this last estimate.
\end{theorem}

\begin{proof}
The reverse triangle inequality and the fact that the positive-part
function is $1$-Lipschitz give
$|w_i(x)-w_i(y)|\le d(x,y)$ and
$|S(x)-S(y)|\le Nd(x,y)$.  The coefficients $w_i(x)/S(x)$ are
nonnegative and sum to one, proving~\eqref{eq:an-glue-bound}.

Write $\alpha=w_i(x)$ and $\beta=w_i(y)$.  If both are positive,
split the difference as
\[
 \alpha f_i(p,x)-\beta f_i(q,y)
 =(\alpha-\beta)f_i(q,y)
   +\alpha\bigl(f_i(p,x)-f_i(q,y)\bigr).
\]
Its norm is at most
$M|\alpha-\beta|+
\alpha(Ld(x,y)+\anE_p+\anE_q)$.
If $\alpha=0$ or $\beta=0$, the same bound remains valid: the
nonzero summand has norm at most $M|\alpha-\beta|$, and the extra
term is nonnegative.  Thus no assertion about an inactive local
value has been used.  Summing gives
\[
 \norm{A_p(x)-A_q(y)}_1
 \le NMd(x,y)+S(x)
       \bigl(Ld(x,y)+\anE_p+\anE_q\bigr).
\]
Finally,
\[
 \frac{A_p(x)}{S(x)}-\frac{A_q(y)}{S(y)}
 =\frac{A_p(x)-A_q(y)}{S(x)}
  +A_q(y)\frac{S(y)-S(x)}{S(x)S(y)}.
\]
Since $\norm{A_q(y)}_1\le MS(y)$ and $S(x)\ge\delta$, the first
term contributes at most $(L+NM/\delta)d(x,y)+\anE_p+\anE_q$,
and the second contributes at most $(NM/\delta)d(x,y)$.
This proves~\eqref{eq:an-glue-lip}.
To prove~\eqref{eq:an-glue-restrict}, subtract $f_j(q,x)$ inside
the convex sum and apply~\eqref{eq:an-overlap} term by term.
\end{proof}

\subsection{Completion, coverage certificates, and exact gluing}

Choose a natural $c\ge2$ with $C\le2^{c-2}$.  For a regular input
name $x$, define the output samples
\begin{equation}
 b_n=U_{n+2}(x_{n+c}).
 \label{eq:an-glue-name}
\end{equation}
Equation~\eqref{eq:an-glue-lip} gives
\[
 \norm{b_m-b_n}_1
 \le C2^{-c}(\anE_m+\anE_n)
      +\anE_{m+2}+\anE_{n+2}
 \le\tfrac12(\anE_m+\anE_n).
\]
Hence the output is a regular name.  Applying the same finite estimate
to equivalent input names proves extensionality.  The induced map on
the completion is $C$-Lipschitz; the proof follows by inserting late
samples and removing the vanishing error.  This argument directly
allows approximate local evaluators, so it does not require an exact
rational value of the local functions.

The denominator hypothesis can itself be certified by a finite net.
Suppose a finite list $v_1,\ldots,v_s$ covers the domain within distance
$\rho$, and a finite calculation verifies
\[
 S(v_j)\ge\delta+N\rho\quad(j<s).
\]
For any covered point choose its supplied covering witness, use the
$N$-Lipschitz bound for $S$, and obtain $S(x)\ge\delta$.  Internally
the coverage witness is a bounded search result or part of the domain
contract.  A claimed arbitrary open cover does not manufacture this
net or its coverage proof.

If the fixed local routines agree on every overlap to all requested
precisions, use~\eqref{eq:an-glue-restrict} with
$\eta=\anE_k$ for arbitrary $k$.  The completed global value equals
each active local value.  At a completed point with positive clearance
inside a chart, sufficiently late rational samples remain in that
chart, so the same equality holds there.  Positive clearance follows
from the robust cover data for at least one chart; exact chart identity
need not be selected by the program.

If one has only a fixed overlap tolerance $\eta$, the theorem gives
approximate gluing with exactly that tolerance.  A different family of
local answers at every precision must additionally have a cross-stage
coherence contract.  Same-stage agreement is insufficient: all charts
could agree on $0$ at even stages and on $1$ at odd stages.  Thus the
construction does not turn independent approximate local solutions
into an exact global solution without a coherence premise.

\paragraph{Named distances and safe support tests.}
If $d(x,c_i)$ is approximated by $\Delta_i(t)$ with error
$\anE_t$, use the inner weight
\[
 \widetilde w_i=\max(0,r_i-\Delta_i(t)-\anE_t).
\]
Then $0\le\widetilde w_i\le w_i$,
$w_i-\widetilde w_i\le2\anE_t$, and
$\widetilde w_i>0$ certifies genuine membership in the local domain.
If $2N\anE_t\le\delta/2$, their sum is at least $\delta/2$.
Replacing the exact weights by these inner weights changes the
normalized sum by at most
\[
 \frac{8MN}{\delta}\anE_t.
\]
Indeed, the numerator changes by at most
$M\sum_i(w_i-\widetilde w_i)$; the denominator correction contributes
the same amount before division by a denominator at least $\delta/2$.
Choose $t$ so this additional error fits the requested output budget.
This supplies a finite safe algorithm even when exact chart membership
is not decidable from a metric name.

The target convexity is essential.  Weighted averages of manifold
points need not lie on the manifold.  A manifold-valued version needs
an explicit retraction from the convex neighborhood, together with its
domain margin and continuity modulus.  None is inferred from the
word ``manifold.''

\section{From approximate solutions to an identified completed point}
\label{sec:residual-recovery}

The repository's coherence obstruction prevents an unrestricted
inference from finite solvability at every horizon to one computable
coherent branch.  A positive theorem must specify the extra information
that rules out that obstruction.  Here it is a modulus bounding the
diameter of the approximate solution sets.  This is the same core
principle already used by \texttt{CoherentMinimizer}; completion now
also identifies the limiting value and proves its uniqueness.

\subsection{The generic residual theorem}

Let $r:D\to\Q_{\ge0}$ be a fixed uniformly continuous residual, with
a supplied dyadic modulus.  Assume a fixed primitive-recursive
procedure $a:\N\to D$ and a primitive-recursive function $\kappa$
satisfy
\begin{align}
 r(a(n))&\le\anE_n,\label{eq:an-residual-witness}\\
 r(u),r(v)\le\anE_{\kappa(k)}
 &\quad\Longrightarrow\quad d(u,v)\le\anE_k.
 \label{eq:an-diameter}
\end{align}
The approximate witnesses are an assumption to prove in the application;
they are not extracted from an assertion that an exact zero exists.
Normalize $\kappa$ to the nondecreasing majorant
\[
 \kappa^*(k)=k+\max_{j\le k}\kappa(j).
\]
It is at least $k$ and at least $\kappa(k)$.

\begin{theorem}[Recovery and identification by a diameter modulus]
\label{thm:an-residual-recovery}
The samples
\[
 z_n=a(\kappa^*(n+2))
\]
form a regular name.  The continuous extension $\anH r$ satisfies
$\anH r(z)=0$.  This is the unique completed zero of $\anH r$.
Every construction is primitive recursive relative to the fixed
displayed procedures.  The same statement holds for a named residual
when the finite witness and diameter contracts include its sample
error allowances.
\end{theorem}

\begin{proof}
For $m\ge n$, both residuals are at most
$\anE_{\kappa^*(n+2)}\le\anE_{\kappa(n+2)}$.
Equation~\eqref{eq:an-diameter} therefore gives
$d(z_m,z_n)\le\anE_{n+2}$, proving regularity.
Also $r(z_n)\le\anE_{n+2}$ because $\kappa^*(n+2)\ge n+2$.

Let $\nu_r$ be the extension modulus furnished by
Theorem~\ref{thm:an-extension}.  To bound $\anH r(z)$ at output
precision $k$, choose
$n\ge\max(k+2,\nu_r(k+2))$.  Density gives
$\anH d(z,\operatorname{const}(z_n))\le\anE_n$, so
$|\anH r(z)-r(z_n)|\le\anE_{k+2}$.
Together with $r(z_n)\le\anE_{k+2}$ this bounds
$|\anH r(z)|$ by $2\anE_{k+2}$ for every $k$.  Hence it is zero.
This is the canonical-evaluation step missing from a Cauchy estimate
alone.

For uniqueness, let $y$ be a regular name with $\anH r(y)=0$.
Fix $k$ and put $t=\kappa^*(k+3)$.  Choose
\[
 n\ge\max(k+3,\nu_r(t+1)),\qquad
 m\ge\max(k+3,t).
\]
Uniform continuity and density imply
$r(y_n)\le\anE_{t+1}$, while
$r(z_m)\le\anE_{m+2}\le\anE_t$.
Both are below $\anE_{\kappa(k+3)}$, so
$d(z_m,y_n)\le\anE_{k+3}$.
The completed distance between $z$ and $y$ is consequently at most
\[
 \anE_m+\anE_{k+3}+\anE_n
 \le3\anE_{k+3}<\anE_k.
\]
This holds at every $k$ and proves $y\sim z$.  Every index in the
proof is an explicitly displayed primitive-recursive composition or
maximum.  No exact solution was chosen at an intermediate stage.
\end{proof}

There are two separate assumptions here.  Approximate solvability
supplies points; the diameter modulus forces any sufficiently accurate
points to agree.  Neither replaces the other.  Mere uniqueness of a
semantic zero does not automatically supply a primitive-recursive
diameter modulus, and a diameter bound without approximate solvability
could be vacuous.

\subsection{A complete contraction instance}

Let $F:D\to D$ be a fixed map with
$d(Fu,Fv)\le qd(u,v)$ for a supplied rational $0\le q<1$.
Choose a starting point $x_0$ and integer $B\ge0$ such that
$d(x_0,Fx_0)\le2^B$.  Let $H\ge1$ be an integer satisfying
$H(1-q)\ge1$.  Both $B$ and $H$ can be found by bounded searches from
the rational data.

For $\gamma=1-q$, the identity
$q(1+\gamma)=1-\gamma^2\le1$ and the finite Bernoulli inequality
$(1+\gamma)^H\ge1+H\gamma\ge2$ imply $q^H\le1/2$.
The Bernoulli inequality is proved by induction on the exponent with
a polynomial inequality as invariant.  Iterating the contraction
gives
\[
 d(F^j x_0,F^{j+1}x_0)\le q^j2^B.
\]
For the residual $r(x)=d(x,Fx)$, set
\[
 a(n)=F^{H(n+B)}x_0.
\]
Then $r(a(n))\le\anE_n$.  This constructs approximate solutions
directly from finite iteration; it assumes no fixed point.

The triangle inequality gives
\[
 (1-q)d(u,v)\le r(u)+r(v).
\]
Choose $c$ with $2^{1-c}\le1-q$.  Then
$\kappa(k)=k+c$ is a diameter modulus.
Moreover $|r(u)-r(v)|\le(1+q)d(u,v)$, so the residual is uniformly
continuous.  Theorem~\ref{thm:an-residual-recovery} now constructs a
unique completed zero.  By the extension theorem, the residual
identifies with $x\mapsto\anH d(x,\anH F(x))$: these two continuous
functions agree on the dense constant names.  Thus the zero is the
unique fixed point of $\anH F$.

The source module \texttt{ContractionIterates} already supplies
quantifier-free induction proofs of contraction estimates for
iterations.  The additional reusable step here is the completed
fixed-point identity and its uniqueness, with the exact dependency
on continuity and density exposed.

\section{A singular inverse that does not need derivative separation}
\label{sec:singular-cube-root}

Derivative separation is a sufficient way to obtain a diameter
modulus, but it is not necessary.  The real cube root is continuous at
zero even though the derivative of $y\mapsto y^3$ vanishes there and
the inverse is not locally Lipschitz.  Its quantitative inverse
modulus is still primitive recursive.

The inspected modules \texttt{StableCubicRoot} and
\texttt{CubicInverse} concern $y^3+y=x$, whose extra linear term
gives inverse Lipschitz constant one.  The example here deliberately
uses $y^3=x$ and includes negative inputs and the singular value zero.
It does not produce a holomorphic branch of a complex cube root near
zero, nor does it decide the sign of a represented real number.

\begin{lemma}[A global cubic separation inequality]
\label{lem:an-cubic-separation}
For all rational $u,v$,
\begin{equation}
 |u-v|^3\le4|u^3-v^3|.
 \label{eq:an-cubic-separation}
\end{equation}
The same inequality holds for represented reals by approximation.
\end{lemma}

\begin{proof}
If $u,v$ have the same sign, change both signs if necessary and assume
$u\ge v\ge0$.  Then
\[
 u^3-v^3=(u-v)(u^2+uv+v^2)\ge(u-v)^3,
\]
because $u^2+uv+v^2-(u-v)^2=3uv\ge0$.
If they have opposite signs, put $s=|u|$ and $t=|v|$.  The exact
identity
\[
 4(s^3+t^3)-(s+t)^3=3(s+t)(s-t)^2\ge0
\]
proves the result.  All case splits compare rational integers, and
all substantive steps are polynomial identities and nonnegative
products.  Passing to real names uses the explicit product error
bounds on a supplied bounded interval and
Lemma~\ref{lem:an-small-error}.
\end{proof}

Let $A_p$ be a fixed regular name of $a\in[-1,1]$.  Semantically
$|A_p-a|\le\anE_p$; internally this is expressed by the
corresponding cross-sample inequalities.  Define
$C_p=\max(-1,\min(1,A_p))$.  Clipping is $1$-Lipschitz and fixes
$[-1,1]$, so $|C_p-a|\le\anE_p$.

For a requested residual index $n$, put
\[
 p=n+4,\qquad J=2^{n+5},\qquad q_j=-1+2j/J\quad(0\le j\le J).
\]
Use bounded search to select the largest $j\le J$ for which
$q_j^3\le C_p$ and return $a_n=q_j$.  There is always at least one
candidate because $q_0=-1$ and $C_p\ge-1$.  All comparisons concern
the rational sample $C_p$, so even the input $a=0$ causes no exact
real comparison.

\begin{proposition}[Primitive-recursive approximate cubic roots]
\label{prop:an-cubic-approx}
The returned point lies in $[-1,1]$ and satisfies
\begin{equation}
 |a_n^3-a|\le\anE_{n+2}.
 \label{eq:an-cubic-residual}
\end{equation}
\end{proposition}

\begin{proof}
If $j=J$, maximality forces $C_p=1$ and the error relative to
$C_p$ is zero.  Otherwise
$q_j^3\le C_p<q_{j+1}^3$.  For $x,y\in[-1,1]$,
\[
 |x^3-y^3|=|x-y||x^2+xy+y^2|\le3|x-y|.
\]
The mesh length is $2/J$, so
$0\le C_p-a_n^3\le6/J$.
Adding the input sample error yields
\[
 |a_n^3-a|\le6\,2^{-(n+5)}+2^{-(n+4)}
 =\tfrac14\,2^{-n}=\anE_{n+2}.
\]
The maximum-search invariant states that the selected index lies
within the finite grid, passes the comparison, and every later
inspected index fails it.  A Boolean fold expresses the last condition;
the invariant is quantifier free.  No intermediate-value theorem or
compactness principle has been used to produce the witness.
\end{proof}

If $|u^3-a|$ and $|v^3-a|$ are at most
$\anE_{3k+3}$, then
\[
 |u-v|^3\le4\bigl(|u^3-a|+|v^3-a|\bigr)
 \le8\anE_{3k+3}=\anE_k^3.
\]
Monotonicity of the cube on nonnegative rationals proves
$|u-v|\le\anE_k$.  Thus
\begin{equation}
 \kappa(k)=3k+3
 \label{eq:an-cubic-kappa}
\end{equation}
is an explicit residual-diameter modulus.  The finite contractions in
this proof remain valid with rational sample errors: use an arbitrary
sufficiently precise common sample $A_t$, apply the rational triangle
inequality, and remove $2\anE_t$ by
Lemma~\ref{lem:an-small-error}.

In particular, the simple schedule
\[
 z_k=a_{3k+9}
\]
already gives a regular name: for $m\ge k$ both residuals meet the
threshold $\kappa(k+2)=3k+9$, so
$|z_m-z_k|\le\anE_{k+2}$.
The residual is continuous on $[-1,1]$, with Lipschitz constant three,
so the recovery proof identifies $z^3=a$ and proves uniqueness.
For a supplied integer bound $B\ge1$ with $|a|\le B$, apply the
construction to $a/B^3\in[-1,1]$ and multiply the output by $B$.
An extra sample shift $s$ with $B\le2^s$ restores the regularity
budget.  This extends the construction to every bounded real-name
family using only rational scaling and bounded doubling.

Moreover, for the outputs $x=\sqrt[3]{a}$ and $y=\sqrt[3]{b}$,
\[
 |x-y|^3\le4|a-b|.
\]
Thus the inverse itself has dyadic modulus $k\mapsto3k+2$.
There is no discontinuous branch decision.  A loss from $k$ output
bits to roughly $3k$ input bits expresses the conditioning at the
singular value.  This resolves the construction at this singularity
without making the much stronger, false assertion that every
singular analytic branch can be chosen continuously.

\section{Formalization obligations for the continuous layer}
\label{sec:analytic-lean-obligations}

The accompanying Lean blueprint names the semantic interfaces and the
finite statements to implement.  These are not completed new
\texttt{Prf PRAAx} proofs.  The reliable order of implementation is
as follows.

\begin{enumerate}
\item Encode rational vectors, their $\ell^1$ norm, the positive-part
operation, and the finite sums in~\eqref{eq:an-glue}.  Denominators
are positive natural codes; signed numerators are handled by the
existing integer arithmetic layer.  Prove the one-summand inequality
in the proof of Theorem~\ref{thm:an-glue} first.

\item Derive the finite weighted-sum bound by induction on the list
length.  The invariant is the cleared rational inequality for the
partial numerator and partial weight sum.  Retain $\delta>0$ as an
explicit premise before cancelling denominators.  This produces an
object theorem for arbitrary finite chart tables and precision
parameters.

\item Formalize the reindexing inequalities for
Theorem~\ref{thm:an-extension} and the diagonal theorem.  Their
internal conclusions are universal closures of rational inequalities
between samples of specified \texttt{PR} programs.  Regularity and
modulus hypotheses are instantiated at the displayed indices.
There is no induction over names or real numbers.

\item Reuse the maximum-normalization and coherence lemmas from
\texttt{CoherentMinimizer}; add canonical-evaluation and uniqueness
lemmas using the completion interface.  The proof goal is the
explicit sample equality formula for the residual name, not merely a
Lean equality in an ambient complete space.

\item Obtain the two cubic identities through the existing polynomial
derivation mechanism.  Derive the same-sign and opposite-sign cases
using decidable rational sign tests.  Implement the finite grid
selector with the existing bounded-search interface; discharge its
endpoint, maximality, mesh, and residual inequalities.  Finally
compose the residual contracts with the explicit cubic diameter
modulus.

\item Only after these deep arithmetic endpoints are available,
state the external semantic theorems about completions, continuous
maps, and real cube roots.  Prove that the represented evaluations
agree with those semantic objects.  An ambient Lean theorem alone
does not replace any of the preceding endpoints.
\end{enumerate}

For a general distance presentation, every object formula additionally
contains the finite distance-query index and its error allowance.
The trust boundary includes the proofs of those distance contracts,
of coverage, and of the local moduli.  These are useful explicit
obligations, rather than implicit appeals to compactness or a
qualitative continuity theorem.

\section{Compact presentations of spaces of functions}\label{sec:functionspaces}

The preceding completion construction describes a point from finite observations. We now make an entire space of continuous functions into a presented sort. This addresses a different difficulty from finding a net in a given finite-dimensional domain: every net center must itself be a valid function. A quantized sample table need not satisfy the original Lipschitz inequalities. Searching for a function extending it would reintroduce an existence problem. The clipped envelope below accepts every table and always produces a valid center.

\subsection{The class and its numerical parameters}

Fix \(d\geq1\), an integer \(M\geq1\), and a nonnegative rational \(L\). Put
\[
 K=[0,1]^d,\qquad d_K(x,y)=\max_{1\leq j\leq d}|x_j-y_j|,
\]
and consider
\[
 \mathcal F_{M,L}=\{f:K\to[-M,M]:
       |f(x)-f(y)|\leq Ld_K(x,y)\}.
\]
Its semantic distance is \(D(f,g)=\sup_{x\in K}|f(x)-g(x)|\). The construction below also defines this supremum as a regular real name, so its use does not presuppose a supremum axiom in \ipra. A rational \(L\) is supplied with a natural number \(s\) satisfying \(L\leq2^s\); one may take a deliberately loose PR bound on \(s\). The case \(L=0\) is included.

Let \(G_m\) be the grid whose coordinates are \(i/2^m\), \(0\leq i\leq2^m\). It has
\[
 s_m=(2^m+1)^d
\]
points and covering radius at most \(h=2^{-m}\). There is a constructive covering procedure: approximate each coordinate of \(x\) by a rational within \(h/4\), clip it into \([0,1]\), and round to the nearest grid coordinate. The total error is at most \(3h/4\). This avoids exact floor or exact comparison on real numbers.

\subsection{A finite envelope is always a legal function}

For any rational table \(q=(q_a)_{a\in G_m}\), define
\begin{equation}
 E_q(x)=\min_{a\in G_m}\bigl(q_a+Ld_K(x,a)\bigr),\qquad
 P_q(x)=\clamp_{[-M,M]}E_q(x).
 \label{eq:envelope}
\end{equation}
Clipping means \(\clamp_{[u,v]}t=\max(u,\min(v,t))\). On rational inputs all these operations are exact PR rational operations. On real inputs min and max are continuous lattice operations; no selector for the exact sign of a real difference is used.

Formula \eqref{eq:envelope} is the finite form of the classical Lipschitz extension envelope \cite{mcshane}. We use it for approximate tables rather than requiring an exactly consistent extension problem.

\begin{lemma}[Envelope bounds]\label{lem:envelope}
For every table \(q\), \(P_q\in\mathcal F_{M,L}\). If \(f\in\mathcal F_{M,L}\), the grid has covering radius \(h\), and
\(\abs{q_a-f(a)}\leq\delta\) for every grid point, then
\[
 f(x)-\delta\ \leq\ E_q(x)\ \leq\ f(x)+\delta+2Lh
\]
for all \(x\in K\), and therefore
\begin{equation}
 D(P_q,f)\leq\delta+2Lh.
 \label{eq:envelope-error}
\end{equation}
\end{lemma}

\begin{proof}
The triangle inequality gives
\(|d_K(x,a)-d_K(y,a)|\leq d_K(x,y)\). Hence each cone \(q_a+Ld_K(-,a)\) is \(L\)-Lipschitz. If two finite families obey \(u_a\leq v_a+c\) for every \(a\), their minima obey \(\min u_a\leq\min v_a+c\). Apply this in both directions to the cones. Clipping is nonexpansive, so it preserves the Lipschitz bound and gives the prescribed range. In particular every table is legal, irrespective of consistency between its entries.

For any \(a\), Lipschitz continuity of \(f\) implies
\[
 q_a+Ld_K(x,a)\geq f(a)-\delta+Ld_K(x,a)\geq f(x)-\delta.
\]
Take the finite minimum. Choose a grid point within \(h\) of \(x\), using the preceding covering procedure. Its cone is bounded above by
\[
 f(a)+\delta+Ld_K(x,a)\leq f(x)+\delta+2Lh.
\]
The minimum is no larger. Since clipping fixes \(f(x)\) and is nonexpansive, it cannot increase the absolute error. All inequalities for rational \(x\), rational samples, and a specified covering witness are finite ordered-rational inequalities. Their bounded conjunctions and minimum scans use the invariants in Theorem~\ref{thm:realization}.
\end{proof}

\subsection{The net and its explicit cardinality}

For precision \(k\), set
\[
 \delta_k=2^{-(k+2)},\qquad m_k=k+3+s,
 \qquad B_k=2M\,2^{k+2}+1.
\]
Allow every entry \(q_a\) to lie in the finite rational set
\[
 Q_k=\{-M,-M+\delta_k,\ldots,M\}.
\]
There are exactly \(B_k^{s_{m_k}}\) tables; repetitions among their envelope functions cause no problem.

\begin{theorem}[A PR net of genuine functions]\label{thm:function-net}
The family
\[
 \mathcal N_k=\{P_q:q\in Q_k^{G_{m_k}}\}
\]
is a \(2^{-k}\)-net in \(\mathcal F_{M,L}\), with the stronger error bound \(2^{-(k+1)}\). Its size is bounded by the explicit PR function
\begin{equation}
 N(k,d,M,s)=\bigl(2M\,2^{k+2}+1\bigr)^{(2^{k+3+s}+1)^d}.
 \label{eq:net-size}
\end{equation}
Its centers have numerical codes and exact PR evaluation at rational points. Given a fixed PR evaluator for \(f\), a suitable center is selected by a PR finite computation. No evaluator for arbitrary PR function codes is required.
\end{theorem}

\begin{proof}
At each grid point query a rational sample \(v_a\) of \(f(a)\) with error at most \(2^{-(k+3)}=\delta_k/2\). Round this rational to the nearest multiple of \(\delta_k\), and clip to \([-M,M]\). Since the interval endpoints are themselves multiples, the result belongs to \(Q_k\), with \(|q_a-f(a)|\leq\delta_k\). This is a rational rounding operation; it is not exact rounding of an unknown real number.

The covering radius is \(2^{-m_k}\), and
\[
 2L\,2^{-m_k}\leq2\,2^s\,2^{-k-3-s}=2^{-(k+2)}=\delta_k.
\]
Lemma~\ref{lem:envelope} gives error at most \(2\delta_k\). Enumerate the tables by their base-\(B_k\) digit words of length \(s_{m_k}\). Computing a digit, a grid point, a cone value, and a finite minimum are all bounded PR operations. Formula~\eqref{eq:net-size} bounds this enumeration; it is very large, but it is primitive recursive even with the displayed numerical parameters variable.

The entire catalog of centers is independent of \(f\). If one wants to locate a center near a supplied \(f\), one asks only the finite list of samples just described. This computation is PR relative to that evaluator, and ordinary PR after fixing a PR evaluator. In contrast, evaluating a center needs only its finite table code. Quantifiers over this function space will use the latter, unconditional evaluation procedure.
\end{proof}

\begin{remark}[All semantic functions versus all function codes]
The semantic covering statement ranges over every bounded \(L\)-Lipschitz function, including functions without a computable name in a classical ambient model. The numerical centers form an explicitly presented dense family. This does not enumerate all those functions or make their pointwise evaluation computable. An extremal value is computable from a supplied evaluator, a dense net, and a uniform continuity modulus; it does not need a listing of every semantic function. This is exactly why the construction supports function quantification without a universal evaluator.
\end{remark}

\subsection{Distance and evaluation are themselves presented maps}

\begin{lemma}[Finite evaluation of the supremum distance]\label{lem:function-distance}
Let \(G\) have covering radius \(h\). Suppose samples \(a_z,b_z\) approximate \(f(z),g(z)\) within \(\eta\), where \(f,g\in\mathcal F_{M,L}\), and put
\[
  Z=\max_{z\in G}|a_z-b_z|.
\]
Then
\begin{equation}
  |Z-D(f,g)|\leq2\eta+2Lh.
 \label{eq:function-distance}
\end{equation}
Moreover evaluation satisfies
\begin{equation}
 |f(x)-g(y)|\leq D(f,g)+Ld_K(x,y).
 \label{eq:evaluation-lip}
\end{equation}
\end{lemma}

\begin{proof}
At each grid point the sampled difference differs from the true absolute difference by at most \(2\eta\). The finite maximum preserves this error. The function \(x\mapsto|f(x)-g(x)|\) is \(2L\)-Lipschitz. Its maximum on a net is no greater than its supremum and is within \(2Lh\) of every value. These two one-sided bounds give \eqref{eq:function-distance}; they also construct the supremum as a Cauchy value when meshes tend to zero. Finally insert \(g(x)\) between \(f(x)\) and \(g(y)\) and use the triangle inequality.
\end{proof}

For two coded envelope centers, their grid values are exact rationals, so \(\eta=0\). Selecting \(h\) from \(k,L\) makes \(Z\) a PR distance evaluator at any desired precision. We therefore do not need an exact decision procedure for their real supremum distance. Triangle and symmetry laws for the distance names follow by the same finite error accounting and are of the form in Definition~\ref{def:presentation}.

For the diameter-at-most-one convention of continuous logic, use the metric \(\overline D=D/(2M)\). The positive scaling factor is known. The same raw nets are still nets at the indicated normalized scale, and evaluation has the explicit bound \(2M\overline D(f,g)+Ld_K(x,y)\). Thus all rescaling and atomic-modulus obligations remain numerical.

To evaluate at a point supplied by a Cauchy name, first request its coordinates at a precision making \(L\) times the input error less than half the requested output error, then evaluate the center on that rational point. For an arbitrary named function one additionally requests an appropriately accurate function value. Equation~\eqref{eq:evaluation-lip} proves both the demand bound and representation independence.

\subsection{Completeness without extracting a subsequence}

\begin{proposition}[Uniform completion]\label{prop:function-complete}
Suppose a sequence \((f_n)\) in \(\mathcal F_{M,L}\) comes with a nondecreasing modulus \(b\) such that
\[
 n,m\geq b(k)\quad\Longrightarrow\quad D(f_n,f_m)\leq2^{-k}.
\]
Then it has a limit in \(\mathcal F_{M,L}\). If the sequence evaluator and \(b\) are fixed PR programs with the indicated finite proofs, the limit evaluator has a PR presentation. Together with Theorem~\ref{thm:function-net}, this supplies completeness and explicit total boundedness of the function-space presentation.
\end{proposition}

\begin{proof}
At input \(x\) and output precision \(k\), evaluate \(f_{b(k+3)}(x)\) with error at most \(2^{-(k+3)}\). When \(k\leq l\), the two chosen function indices are at least \(b(k+3)\), so the two returned values differ by at most
\[
 2^{-(k+3)}+2^{-(k+3)}+2^{-(l+3)}.
\]
This is a regular Cauchy bound with slack. It defines \(f(x)\) directly; there is no choice of a limit point independently for each \(x\). The same estimate, with \(l\) arbitrarily large, gives uniform convergence. The inequalities \(|f_n(x)|\leq M\) and \(|f_n(x)-f_n(y)|\leq Ld_K(x,y)\) pass through the explicit error estimates to give the corresponding inequalities for \(f\). At the arithmetic level the conclusion is the all-precision family of rational inequalities, not an induction on a statement quantified over all functions.

If necessary, replace a supplied modulus by \(\max_{j\leq k}b(j)\). This is a PR normalization. The evaluator just described is a fixed composition of the sequence evaluator with \(b\) and the requested precision. If the sequence was oracle-relative, so is this construction until the input programs are fixed.
\end{proof}

``Compact'' here means this explicit complete, totally bounded presentation. We do not extract a convergent subsequence from an arbitrary sequence. A real-valued functional supplied with a PR evaluator and a PR uniform-continuity modulus has a supremum value and PR approximate maximizers on these nets, but an exact maximizing point needs additional information such as the residual-diameter condition. These distinctions agree with the constructive treatment of compactness \cite{bishop}.

\subsection{A quantitative function-object correspondence}

Let \(A\) be another presented metric space. There is a precise useful form of currying that does not pretend that all metric spaces form a cartesian closed category.

\begin{proposition}[Evaluation and currying with fixed bounds]\label{prop:currying}
Suppose \(H:A\times K\to[-M,M]\) satisfies
\begin{equation}
 |H(a,x)-H(b,y)|\leq C d_A(a,b)+L d_K(x,y).
 \label{eq:joint-budget}
\end{equation}
Then \(\widehat H(a)(x)=H(a,x)\) defines a \(C\)-Lipschitz map \(A\to\mathcal F_{M,L}\). Conversely, a \(C\)-Lipschitz map into \(\mathcal F_{M,L}\) defines \(H\) satisfying \eqref{eq:joint-budget}. The two transformations are inverse on extensional maps and have the finite evaluation procedure given above.
\end{proposition}

\begin{proof}
Set \(a=b\) to see that every section belongs to the function class. Set \(x=y\) and take the supremum over \(x\) to get
\(D(\widehat H(a),\widehat H(b))\leq C d_A(a,b)\). Conversely use \eqref{eq:evaluation-lip} and the assumed function-space distance bound. Both composite transformations return the original value at every pair \((a,x)\), proving the inverse equations. The pointwise finite-grid estimates prove the same equations to every requested precision; closure follows from fixed PR composition and the supplied moduli.
\end{proof}

This is a bounded function object: the bounds \(M,L,C\) are part of the interface. The completion-extension theorem provides its categorical transport, while the continuous model semantics can treat it as a compact sort. Quantifiers over its elements are implemented by finite table codes for the envelopes, not by quantifying over programs for arbitrary functions.

\subsection{Why the regularity bounds cannot simply be removed}

The unit ball of all continuous functions on \([0,1]\) has no finite supremum-metric net at every precision. Indeed let
\(f_n(x)=\min(1,2^n x)\). For \(m>n\), at \(x=2^{-m}\) the two values are \(1\) and at most \(1/2\). Thus \(D(f_n,f_m)\geq1/2\). A ball of radius less than \(1/4\) contains at most one such function, so no finite collection covers the family. Each example has an elementary PR evaluator. The failure is uniform regularity, not computability of the individual members. This simple obstruction explains why compact function sorts must carry shared range and continuity data.

% Continuous model-theoretic chapter. Shared macros supplied by the main paper.
\providecommand{\val}[1]{\lVert #1\rVert}
\providecommand{\qval}{\operatorname{qval}}
\providecommand{\Amp}{\operatorname{Amp}}
\providecommand{\CM}{\mathcal M}
\providecommand{\CN}{\mathcal N}
\section{Continuous model theory from finite correspondences}
\label{sec:continuous-models}

The repository already contains presented categories, a presentation-wise
Yoneda theorem, finite limits and colimits, and a setoid interpretation of
intuitionistic first-order derivations. Repeating those constructions would
not solve the present analytic problem. The contribution of this section is
a quantitative interpretation of \emph{real-valued formulas on continuous
structures}. Its central estimate is unusually economical: a two-sided
correspondence with atomic error $\eta$ preserves every formula built from
$\min,\max,1-(-),\inf,\sup$ with the \emph{same} error $\eta$.
Quantifier alternation does not multiply the error.

This is a restricted, effective part of continuous model theory. In the usual
semantic setting, predicates take values in $[0,1]$ and quantifiers are
infima and suprema; uniform continuity makes passage to completions
meaningful. These conventions and their general model-theoretic setting
are described in Ben Yaacov--Berenstein--Henson--Usvyatsov,
\emph{Model theory for metric structures}, Sections 2--4
(\url{https://math.univ-lyon1.fr/~begnac/articles/mtfms.pdf}).
Here we supply finite algorithms and arithmetic error proofs for specified
compact presentations. We do not import that work's compactness,
ultraproduct, saturation, or model-existence theorems into iPRA.

\subsection{The syntax and the scope of a presentation}

Fix a finite collection of sorts and a finite collection of atomic
recipes. An atomic recipe $P$ has a specified finite list of argument sorts.
For a semantic interpretation its value is a uniformly continuous function
into $[0,1]$. A formula is obtained from atomic recipes and rational constants
in $[0,1]$ by the operations
\begin{equation}
 \min(\varphi,\psi),\qquad
 \max(\varphi,\psi),\qquad
 1-\varphi,\qquad
 \inf_{x:s}\varphi,\qquad
 \sup_{x:s}\varphi .
 \label{eq:cm-syntax}
\end{equation}
Free variables carry sorts. Quantification always uses a nonempty sort.
The operations in \eqref{eq:cm-syntax} are numerical operations on formula
values, not replacements for the intuitionistic implication and
quantifiers of the ambient arithmetic.

The fixed-atom convention includes ordinary terms. For a fixed formula,
an occurrence $P(t_1(\bar x),\ldots,t_r(\bar x))$ can first be compiled into
one atomic recipe, with its own approximation procedure and modulus.
The syntax tree of each $t_i$ is fixed in advance. This avoids requiring a
finite grid to be closed under the structure's operations. Such a grid
usually is not a substructure: rounded addition on a grid, for example,
need not satisfy the original algebraic identities.
The approximation evaluates the original composite predicates at grid
points. It does not replace the operations by arbitrary rounded tables.

For clarity we first use the basis \eqref{eq:cm-syntax}.
One may also use a specified rational-evaluable connective
$C:[0,1]^r\to[0,1]$ with a proved rational Lipschitz bound $L_C$ for the
maximum metric. Define its formula amplification by
\begin{equation}
\begin{split}
 A_P&=1,\qquad A_c=0,\\
 A_{\min(\varphi,\psi)}=A_{\max(\varphi,\psi)}
   &=\max(A_\varphi,A_\psi),\qquad A_{1-\varphi}=A_\varphi,\\
 A_{\inf_x\varphi}=A_{\sup_x\varphi}&=A_\varphi,\\
 A_{C(\varphi_1,\ldots,\varphi_r)}
   &=L_C\max_i A_{\varphi_i}.
\end{split}
\label{eq:cm-amplification}
\end{equation}
Use $A=\max(1,A_\varphi)$ when choosing a precision budget.
For the core basis, $A_\varphi\leq1$ regardless of formula size or
quantifier depth. Truncated addition is permitted by taking $L_C=2$;
averaging has $L_C=1$. A separate proof or modulus is needed for any
additional connective. Arbitrary continuous connectives are not silently
treated as primitive recursive.

\subsection{The sharp finite transfer theorem}

Let $A_s$ and $B_s$ be nonempty finite sets for each sort $s$. Atomic
interpretations are rational-valued tables. For finite structures,
$\inf$ and $\sup$ mean the minimum and maximum of a finite nonempty list.
A \emph{total correspondence} consists of relations
$R_s\subseteq A_s\times B_s$ such that every member of either carrier
has a related member of the other. Assume
\begin{equation}
 \bigl|P^A(a_1,\ldots,a_r)-P^B(b_1,\ldots,b_r)\bigr|
 \leq\eta
 \quad\text{whenever }a_iR_{s_i}b_i\text{ for every }i.
 \label{eq:cm-atomic-distortion}
\end{equation}
The hypothesis concerns \emph{all related tuples}, not merely tuples
chosen by one matching algorithm. This distinction matters when a variable
occurs more than once or several variables are quantified.

\begin{theorem}[Finite approximate elementarity]
\label{thm:cm-transfer}
For every fixed formula $\varphi(\bar x)$ and coordinatewise related
assignments $\bar a,\bar b$,
\begin{equation}
 \bigl|\varphi^A(\bar a)-\varphi^B(\bar b)\bigr|
 \leq A_\varphi\eta.
 \label{eq:cm-transfer}
\end{equation}
In particular the bound is $\eta$ for the basis
\eqref{eq:cm-syntax}. The finite correspondence and all numerical tests
can be coded by natural numbers. For each fixed formula this theorem
has an iPRA proof using rational arithmetic, bounded scans, and
quantifier-free induction.
\end{theorem}

\begin{proof}
For two pairs of numbers satisfying $|u-u'|,|v-v'|\leq e$, the inequalities
$u\leq u'+e$ and $v\leq v'+e$ imply
$\max(u,v)\leq\max(u',v')+e$. Interchanging the pairs proves the
reverse inequality. The same argument, or complementation, proves the
minimum case. Also
$|(1-u)-(1-u')|=|u-u'|$. A connective $C$ obeys the asserted estimate by
its Lipschitz hypothesis.

It remains to prove the quantifier step without introducing an extra
error. Write $u(a)$ and $v(b)$ for the two subformula values after fixing
the already related parameters. Suppose $|u(a)-v(b)|\leq e$ whenever
$aRb$. Choose an index $a_*$ attaining $\max_Au$ by a finite scan.
Totality gives $b$ with $a_*Rb$, and therefore
\[
 \max_Au=u(a_*)\leq v(b)+e\leq\max_Bv+e.
\]
Choose an index attaining $\max_Bv$ and use the other half of totality
to obtain the reverse inequality. Thus the maxima differ by at most
$e$. For minima, choose $a_*$ attaining $\min_Au$ and a related $b$:
$\min_Bv\leq v(b)\leq u(a_*)+e=\min_Au+e$.
The reverse inequality follows from a minimizer on $B$.
This proves both quantifier cases with no increase in $e$.
Structural induction on the fixed syntax proves
\eqref{eq:cm-transfer}.

For the arithmetic implementation, represent a finite set by its length
and entries, a relation by a Boolean matrix, and rational values by
signed numerator and positive denominator. For the core basis a common
denominator reduces the proof to inequalities on natural numerators.
A finite scan stores its current extremal value and the first index
attaining it. Its invariant says that the index is in the scanned prefix,
its value equals the current extremum, and every scanned entry lies on
the appropriate side of that value. All these assertions are bounded.
After Booleanization they form a quantifier-free induction formula.
The uses of totality eliminate bounded existential witnesses; they do
not choose an infinite family.
Formula induction is a metatheoretic construction of one iPRA proof per
fixed syntax tree. No induction on an unbounded semantic satisfaction
predicate is asserted.
\end{proof}

Two-sided totality is essential. A one-point subset of $[0,1]$ preserves
the distance predicate on its own points exactly, but its value of
$\sup_x\sup_y d(x,y)$ is zero, whereas the interval's value is one.
The missing part of the comparison is coverage of the interval.

\begin{proposition}[Composition of approximate interpretations]
\label{prop:cm-compose}
If $R$ is a total correspondence of atomic error $\eta$ from $A$ to $B$
and $S$ one of error $\theta$ from $B$ to $C$, their relational composite
is total and has atomic error at most $\eta+\theta$.
\end{proposition}
\begin{proof}
For $a$ choose a related $b$ and then a related $c$; the reverse direction
uses the other two totality assertions. For a tuple related by the
composite, take the finitely many intermediate entries $b_i$ witnessing
its coordinates. Applying the two atomic inequalities and adding them
gives the stated bound. Finite relational composition is a bounded
disjunction over the middle carrier, so its table is primitive recursive.
Associativity follows by reassociating the two finite existential
quantifiers; the diagonal relation is an identity. Reversing a
correspondence preserves its error.
\end{proof}

Consequently finite presentations and total correspondences form a
category whose errors compose additively. The value of each sentence
is stable under its arrows with the amplification in
\eqref{eq:cm-amplification}. This supplies an abstract comparison
calculus for continuous observations. It does not assert that every
correspondence is the graph of a continuous map, or that arbitrary
metric spaces have been assembled into an internal category of iPRA.

\subsection{Primitive-recursive compact presentations}
\label{subsec:cm-presentations}

For each sort $s$, take a rational pseudometric $d_s$ on numerical
codes $D_s$ with diameter at most one, and a fixed primitive-recursive
procedure producing a nonempty finite list $N_s(n)$ of codes. If a
validity predicate is needed, all operations are totalized outside it,
and every assertion below is guarded by validity. The central net
contract is
\begin{equation}
 \forall a\in D_s\ \exists j<|N_s(n)|\quad
 d_s(a,N_s(n)_j)\leq 2^{-n}.
 \label{eq:cm-net}
\end{equation}
Its proof is supplied in iPRA; it is not inferred by inspecting a
finite list. Because distances are rational, a least suitable index
can alternatively be found by a bounded scan.

For each fixed atomic recipe $P$, supply rational approximants
$p_P(\bar a,p)\in[0,1]$ and a nondecreasing primitive-recursive
modulus $\mu_P$. Their arithmetical contracts are coherence and stability:
\begin{equation}
\begin{split}
 |p_P(\bar a,p)-p_P(\bar a,q)|
   &\leq 2^{-p}+2^{-q},\\
 d_\infty(\bar a,\bar b)\leq2^{-\mu_P(h)}
 \quad&\Longrightarrow\quad
 |p_P(\bar a,p)-p_P(\bar b,q)|
   \leq2^{-p}+2^{-q}+2^{-h}.
\end{split}
\label{eq:cm-atom-contract}
\end{equation}
Here $d_\infty$ is the maximum of the coordinate distances, and the
maximum over an empty tuple is zero. These are first-order statements
about fixed numerical functions. Denominators are cleared before
placing them in the object language. One can require dyadic
approximants, with a fixed extra factor for the formula's rational
constants, to make the common-denominator implementation immediate.

The usual semantic description follows constructively: the coherent
sequence $p_P(\bar a,p)$ is a real number; its value is uniformly
continuous on the dense presentation and hence extends to the
completion. This interpretation takes place outside first-order iPRA.
It is not an enumeration of all real numbers by primitive-recursive
program codes.

Let
\begin{equation}
 U(h)=\max\bigl(h,\mu_P(j):
       P\text{ occurs in the fixed formula},\ j\leq h\bigr).
 \label{eq:cm-common-modulus}
\end{equation}
This is a primitive-recursive, nondecreasing common modulus. The
additional maximum over $j\leq h$ makes monotonicity explicit even if
the original moduli were not normalized.

Given dense free parameters $\bar a$, append their entries to the
appropriate lists $N_s(n)$, obtaining finite nonempty carriers
$K_s(n,\bar a)$. Duplicates cause no problem. Interpret every atom
on these carriers by $p_P(-,p)$ and evaluate the formula by finite
minimum and maximum. Denote the resulting rational number by
\begin{equation}
 E_\varphi(\bar a;n,p).
 \label{eq:cm-evaluator}
\end{equation}
The same carrier is used at every quantifier in a single evaluation.
Appending free parameters permits the evaluation to be described
literally as a finite structure with an assignment. Finite supports
may contain any additional specified points: none of the estimates
depends on which extra points were appended.

\begin{theorem}[Cross-mesh estimate]
\label{thm:cm-cross-mesh}
If $n,m\geq U(h)$, then
\begin{equation}
 |E_\varphi(\bar a;n,p)-E_\varphi(\bar a;m,q)|
 \leq A_\varphi\bigl(2^{-p}+2^{-q}+2^{-h}\bigr).
 \label{eq:cm-cross-mesh}
\end{equation}
More generally the same estimate compares parameters $\bar a,\bar b$
when $d_\infty(\bar a,\bar b)\leq2^{-U(h)}$.
\end{theorem}
\begin{proof}
Relate an entry of the first carrier to an entry of the second whenever
their distance is at most $2^{-\min(n,m)}$, in the equal-parameter
case. Every entry has a partner by \eqref{eq:cm-net}: the net contained
in the other carrier covers \emph{all} dense codes, including any
appended parameters. Thus the relation is total both ways.
Identical free parameters are related.
The inequality $n,m\geq U(h)$ and
\eqref{eq:cm-atom-contract} give atomic distortion at most
$2^{-p}+2^{-q}+2^{-h}$ on every related tuple. Apply
Theorem~\ref{thm:cm-transfer}.

For different parameters, use instead the relation
$d_s(a,b)\leq2^{-U(h)}$. Net coverage again gives totality; the specified
parameter pairs now belong to this relation by hypothesis. The same
atomic bound and the same transfer theorem apply.
Every relation test and carrier search in this proof is finite and
rational.
\end{proof}

Fix an integer $\ell$ with $A=\max(1,A_\varphi)\leq2^\ell$ and define
\begin{equation}
 p_k=k+\ell+2,\qquad n_k=U(p_k)+1,\qquad
 v_\varphi(\bar a,k)=E_\varphi(\bar a;n_k,p_k).
 \label{eq:cm-canonical-value}
\end{equation}
Then $v_\varphi$ is primitive recursive. Indeed the syntax is fixed,
each atomic procedure and net generator is fixed, and every quantifier
is a bounded fold over a primitive-recursively generated list.
Nested quantifiers can make this computation extremely large; no
polynomial or elementary complexity bound is claimed.

\begin{corollary}[A real-valued interpretation with an iPRA Cauchy proof]
\label{cor:cm-cauchy}
For all dense parameters and all $k,l$,
\begin{equation}
 |v_\varphi(\bar a,k)-v_\varphi(\bar a,l)|
 \leq \tfrac12\bigl(2^{-k}+2^{-l}\bigr).
 \label{eq:cm-cauchy}
\end{equation}
For each fixed formula, the contracts
\eqref{eq:cm-net}--\eqref{eq:cm-atom-contract} therefore yield an iPRA
derivation of the regular-name condition for the fixed numerical
function $v_\varphi$.
\end{corollary}
\begin{proof}
In \eqref{eq:cm-cross-mesh} set $h=\min(p_k,p_l)$.
Monotonicity of $U$ makes its mesh hypotheses hold. Since
$2^{-h}=\max(2^{-p_k},2^{-p_l})\leq2^{-p_k}+2^{-p_l}$, the right side
is at most
$2A(2^{-p_k}+2^{-p_l})
\leq(2^{-k}+2^{-l})/2$.
This is a rational calculation for arbitrary numerical $k,l$.
No limit is taken inside the object theory.
\end{proof}

Define $\varphi^\CM(\bar a)$ externally to be the real represented by
this regular name. The comparison theorem gives two useful consequences:
\begin{equation}
\begin{aligned}
 |E_\varphi(\bar a;n,p)-\varphi^\CM(\bar a)|
   &\leq A_\varphi(2^{-p}+2^{-h})
       &&(n\geq U(h)),\\
 |\varphi^\CM(\bar a)-\varphi^\CM(\bar b)|
   &\leq A_\varphi2^{-h}
       &&(d_\infty(\bar a,\bar b)\leq2^{-U(h)}).
\end{aligned}
\label{eq:cm-semantic-bounds}
\end{equation}
For the first inequality, compare the given evaluation with
\eqref{eq:cm-canonical-value} at successively finer precisions using
\eqref{eq:cm-cross-mesh}, with $h$ fixed. The term $2^{-q}$ tends to
zero and all remaining terms are fixed.
The second follows similarly from the different-parameter estimate.
These are constructive properties of the represented real values.
They are semantic adequacy statements, separate from the finite
arithmetical theorem \eqref{eq:cm-cross-mesh}.

Thus $\mu_\varphi(k)=U(k+\ell)$ is a modulus on the dense presentation.
The completion-extension theorem, Theorem~\ref{thm:an-extension},
gives the safe completed-parameter modulus
\begin{equation}
 \nu_\varphi(k)=\mu_\varphi(k+2)+2=U(k+\ell+2)+2.
 \label{eq:cm-completed-modulus}
\end{equation}
The margin matters because a closed-threshold modulus on dense points
need not extend with its original indices.
For a tuple $\bar x$ of regular names, whose $j$th entries are
$\bar x_j$, an explicit output schedule is
\begin{equation}
 j_k=U(k+\ell+3)+2,\qquad p=k+\ell+2,\qquad n=U(p)+1,
 \qquad E_\varphi(\bar x_{j_k};n,p).
 \label{eq:cm-named-parameters}
\end{equation}
The distance from $\bar x_{j_k}$ to its completed tuple is at most
$2^{-j_k}$, so \eqref{eq:cm-completed-modulus} bounds the parameter
error by $2^{-(k+1)}$. The first inequality of
\eqref{eq:cm-semantic-bounds}, with $h=p$, bounds the evaluation
error by $2A_\varphi2^{-p}\leq2^{-(k+1)}$.
Thus the returned rational has total error at most $2^{-k}$.
This finite schedule is primitive recursive after fixing the input
name routines.
For arbitrary external names this is a finite-query construction
relative to the names; it is not a new first-order quantification over
all such names.

\subsection{The extension to distances given by regular names}
\label{subsec:cm-named-distance}

The exact rational pseudometric was a simplifying presentation choice.
The function sort in Section~\ref{sec:functionspaces} has a
primitive-recursive \emph{approximation} to its supremum distance;
the preceding theorem must therefore be extended before it is applied
to that sort. The required extension preserves a decidable comparison
relation by reserving an explicit margin.

Let $\Delta_s(a,b,r)$ be a symmetric rational distance sampler,
with values differing from the represented distance by at most $2^{-r}$.
Symmetry can be imposed by averaging the samplers at $(a,b)$ and $(b,a)$.
Its arithmetic coherence theorem is
\begin{equation}
 |\Delta_s(a,b,r)-\Delta_s(a,b,t)|\leq2^{-r}+2^{-t}.
 \label{eq:cm-distance-coherence}
\end{equation}
Define the first-order upper-bound formula
\begin{equation}
 \operatorname{Upper}_s(a,b,\rho)\quad\equiv\quad
 \forall r\ 
    \Delta_s(a,b,r)\leq \rho+2^{-r}.
 \label{eq:cm-upper}
\end{equation}
This formula is not decidable and will not be used as a scan condition.
Its semantic meaning is exactly $d_s(a,b)\leq\rho$.
Supply the net contract
$\forall a\ \exists j<|N_s(n)|\
\operatorname{Upper}_s(a,N_s(n)_j,2^{-n})$
and the atomic stability contract \eqref{eq:cm-atom-contract} with
the metric premise replaced by the conjunction of these
$\operatorname{Upper}$ formulas at radius $2^{-\mu_P(h)}$.
These are explicit object-theory obligations for the presentation.

\begin{lemma}[A decidable relation for named distances]
\label{lem:cm-named-relation}
Fix $h$, put $u=U(h)$ and $\delta=2^{-u}$, and take nets of indices
$n,m\geq u+2$. On their augmented finite carriers define
\begin{equation}
 aR_sb\quad\Longleftrightarrow\quad
 \Delta_s(a,b,u+3)\leq \tfrac58\delta.
 \label{eq:cm-slack-relation}
\end{equation}
Then $R_s$ is total both ways, identical free parameters are related,
and every related pair satisfies
$\operatorname{Upper}_s(a,b,\tfrac34\delta)$.
Consequently the cross-mesh estimate
\eqref{eq:cm-cross-mesh} holds for named distances under the stronger
mesh assumptions $n,m\geq U(h)+2$.
\end{lemma}
\begin{proof}
The other net covers a specified carrier point within
$2^{-m}\leq\delta/4$. Instantiate its $\operatorname{Upper}$
certificate at precision $u+3$:
the sampled distance to the covering center is at most
$\delta/4+\delta/8=3\delta/8$, and therefore passes
\eqref{eq:cm-slack-relation}. This proves each coverage direction;
symmetry handles orientation. A diagonal pair has sampled distance
at most $\delta/8$.

For a pair passing the test, coherence at an arbitrary precision $r$
gives
\[
 \Delta_s(a,b,r)
 \leq \Delta_s(a,b,u+3)+2^{-(u+3)}+2^{-r}
 \leq \tfrac34\delta+2^{-r}.
\]
Generalizing over $r$ proves the $\operatorname{Upper}$ assertion.
Since $\tfrac34\delta\leq2^{-\mu_P(h)}$, the supplied atomic stability
contract gives the atomic distortion used in
Theorem~\ref{thm:cm-transfer}. The relation test itself is rational
and primitive recursive. Although an all-precision upper bound is
derived, no induction formula contains that unbounded quantifier:
$r$ was arbitrary in one arithmetic inequality.
\end{proof}

In this version replace \eqref{eq:cm-canonical-value} by
\begin{equation}
 p_k=k+\ell+2,\qquad
 n_k=U(p_k)+2,\qquad
 v_\varphi(\bar a,k)=E_\varphi(\bar a;n_k,p_k).
 \label{eq:cm-named-schedule}
\end{equation}
The Cauchy calculation is unchanged. To compare different parameters,
a distance at most $2^{-(U(h)+1)}=\delta/2$ suffices:
its sampled distance is at most $5\delta/8$, so those parameter pairs
also pass the relation test. The resulting dense-parameter modulus is
$U(k+\ell)+1$. The safe modulus after completion is consequently
\begin{equation}
 \nu_\varphi^{\mathrm{named}}(k)=U(k+\ell+2)+3.
 \label{eq:cm-completed-named-modulus}
\end{equation}
For free parameter names, replace the schedule
\eqref{eq:cm-named-parameters} by
$j_k=U(k+\ell+3)+3$, $p=k+\ell+2$, and $n=U(p)+2$.
The same two errors, each at most $2^{-(k+1)}$, prove output
accuracy $2^{-k}$.
The first inequality of \eqref{eq:cm-semantic-bounds} remains valid
with $n\geq U(h)+2$.

The quantifier-adequacy proof below uses only this version of the
cross-mesh estimate, its uniform parameter modulus, and finite net
coverage. In the named-distance case take $n=U(p)+2$ throughout
that proof: its net radius is at most $2^{-(U(p)+2)}$, which is
smaller than the parameter radius just specified. Thus the constants
in its displayed value and witness estimates are unchanged.
In particular, the rational supremum-distance sampler of
Lemma~\ref{lem:function-distance} and the envelope net satisfy exactly
the numerical interfaces required to use a compact function class as
a quantified sort. No exact comparison of two real distances is used.

\subsection{Why the interpretation really has infimum and supremum}
\label{subsec:cm-quantifier-adequacy}

The construction above did not assume that every bounded subset of the
reals has a supremum. We now prove that its numerical quantifier
operations have exactly the intended meanings.

\begin{theorem}[Constructive quantifier adequacy]
\label{thm:cm-adequacy}
The values just constructed respect every connective in
\eqref{eq:cm-syntax} and satisfy
\begin{equation}
 (\sup_{y:s}\psi)^\CM(\bar a)
    =\sup_{b\in\widehat D_s}\psi^\CM(\bar a,b),
 \qquad
 (\inf_{y:s}\psi)^\CM(\bar a)
    =\inf_{b\in\widehat D_s}\psi^\CM(\bar a,b).
 \label{eq:cm-adequacy}
\end{equation}
Here the displayed suprema and infima exist because the proof constructs
them. In addition, a primitive-recursive bounded scan supplies a dense
point whose value is within $2^{-k}$ of the requested extremum at
precision $k$.
\end{theorem}
\begin{proof}
For a connective, use one common sufficiently fine mesh and atomic
precision for its immediate subformulas. Their finite evaluations obey
the connective equation exactly. The uniform estimates
\eqref{eq:cm-semantic-bounds}, together with the connective's Lipschitz
bound, pass that equation to the represented values. It is immaterial
that the independently defined canonical names might use different
precisions: \eqref{eq:cm-cross-mesh} compares them.

Consider a supremum and write $S$ for its constructed value. At mesh
$n=U(p)+1$, let
\[
 M_{n,p}=\max_{b\in K_s(n,\bar a)}
                 E_\psi(\bar a,b;n,p),
\]
where all subformula quantifiers use the fixed carriers of this
evaluation. The extra-support version of
\eqref{eq:cm-semantic-bounds} yields, uniformly in these $b$,
\begin{equation}
 |E_\psi(\bar a,b;n,p)-\psi^\CM(\bar a,b)|
       \leq 2A_\psi2^{-p}.
 \label{eq:cm-child-error}
\end{equation}
Given any dense $c$, choose $b$ in the finite net with
$d_s(c,b)\leq2^{-n}\leq2^{-U(p)}$.
The second inequality of \eqref{eq:cm-semantic-bounds} and then
\eqref{eq:cm-child-error} give
\[
 \psi^\CM(\bar a,c)
 \leq\psi^\CM(\bar a,b)+A_\psi2^{-p}
 \leq M_{n,p}+3A_\psi2^{-p}.
\]
Taking arbitrarily fine $p$ proves $\psi^\CM(\bar a,c)\leq S$.
Uniform continuity extends this upper-bound property to every completed
point $c$.

Conversely, suppose $z$ bounds $\psi^\CM(\bar a,b)$ above at every
dense point. For every member of the finite net,
\eqref{eq:cm-child-error} gives
$E_\psi(\bar a,b;n,p)\leq z+2A_\psi2^{-p}$.
Taking the finite maximum and then the represented limit gives
$S\leq z$. Hence $S$ is the least upper bound. This proves existence
and the supremum equation without selecting an infinite sequence of
maximizers. The infimum proof reverses all inequalities.

Finally choose by a finite scan a $b_*$ attaining $M_{n,p}$.
The first estimate in \eqref{eq:cm-semantic-bounds}, applied to the
supremum formula, bounds $|S-M_{n,p}|$ by $2A_\psi2^{-p}$.
Together with \eqref{eq:cm-child-error},
\[
 S-\psi^\CM(\bar a,b_*)\leq4A_\psi2^{-p}.
\]
Choose $p=k+\ell+2$ with $A_\psi\leq2^\ell$. The right side is at most
$2^{-k}$. The maximizing index, its point code, and all precision
choices are primitive recursive. The minimizing case is identical.
\end{proof}

There need be no coherent choice of the approximate witnesses as $k$
varies. The theorem constructs a real value and finite witnesses for
each requested accuracy. Constructing an \emph{exact} optimizing point
is a further problem and may need a uniqueness or separation modulus.
This distinction is exactly what prevents the argument from
reintroducing an unproved branch-selection principle.

\subsection{A worked sentence with genuine quantifier alternation}

Take the sort $[0,1]$, rational dense points, dyadic grids, and the
atomic predicate $P(x,y)=|x-y|$. For this atom,
\[
 \bigl||x-y|-|x'-y'|\bigr|
 \leq |x-x'|+|y-y'|
 \leq2d_\infty((x,y),(x',y')),
\]
so $\mu_P(h)=h+1$ is valid. Consider
\begin{equation}
 \Phi=\inf_x\sup_y |x-y|.
 \label{eq:cm-radius-sentence}
\end{equation}
For a fixed $x$, the triangle inequality on the line shows
$|x-y|\leq\max(x,1-x)$ for $0\leq y\leq1$, while $y=0,1$
give equality. Also $\max(x,1-x)\geq1/2$ because one of $x,1-x$
is at least their average; at $x=1/2$ the maximum is $1/2$.
Consequently $\Phi=1/2$.

Every dyadic grid containing $0,1,1/2$ already returns this exact value
under the finite evaluator. More complicated alternating expressions
need not stabilize exactly, but the same cross-mesh estimate applies.
The logical construction handles $\inf\sup\inf\sup$ without allocating
an additional precision loss for each quantifier: the common net
correspondence handles all of them simultaneously. The price of more
quantifiers is the number of bounded evaluations, not a deterioration
of the proven error bound.

The sorts need not consist of scalar points. In particular, using
the named-distance extension in
Section~\ref{subsec:cm-named-distance}, the compact
class of bounded $L$-Lipschitz functions constructed elsewhere in this
report can be a sort. Evaluation $P(f,x)=f(x)$, after affine rescaling
to $[0,1]$, is uniformly continuous because
\[
 |f(x)-g(y)|\leq\|f-g\|_\infty+L\|x-y\|_\infty.
\]
The supplied finite nets consist of actual Lipschitz functions, so a
formula may quantify over that function class and over points, using
the same theorem. This is a controlled interpretation of apparently
higher-order analytic statements. It is restricted to the specified
compact equicontinuous class; it does not quantify over all continuous
functions or over all primitive-recursive programs.

\subsection{A witness form of the Tarski--Vaught argument}

Global compactness can be replaced by more local data when only a
specified family of formulas is of interest. Let $A$ be a nonempty
subdomain of a structure $M$, with all atomic interpretations obtained
by restriction. Fix a finite family of formulas closed under
subformulas. For every subformula $\inf_y\psi(\bar x,y)$, suppose a
specified finite nonempty list $W_{\psi,\bar a,k}\subseteq A$ satisfies
\begin{equation}
 \forall b\in M\ \exists w\in W_{\psi,\bar a,k}\quad
 \psi^M(\bar a,w)\leq\psi^M(\bar a,b)+2^{-k}.
 \label{eq:cm-witness-test}
\end{equation}
For a supremum use the reverse inequality. The list generator and
the proof of its contract are input data. In an arithmetic application
they must have numerical realizations and iPRA proofs; an informal
existence assertion is not a generator.

\begin{proposition}[Supplied witnesses yield elementary observations]
\label{prop:cm-tarski-vaught}
For the specified subformula family, the values computed over $A$
and over $M$ agree at every tuple from $A$, provided those values
are constructed. With a uniform witness tolerance $\tau$ instead of
arbitrarily small $2^{-k}$, the error for the core basis is at most
$q(\varphi)\tau$, where $q$ is defined by
\begin{equation}
\begin{aligned}
 q(P)=q(c)&=0,\qquad q(1-\varphi)=q(\varphi),\\
 q(\min/\max(\varphi,\psi))&=\max(q(\varphi),q(\psi)),\\
 q(\inf/\sup_y\varphi)&=q(\varphi)+1.
\end{aligned}
 \label{eq:cm-rank-budget}
\end{equation}
\end{proposition}
\begin{proof}
The atomic case is equality, and the connective steps follow from
nonexpansiveness. Suppose the subformula values differ by at most
$e$ on tuples from $A$, and consider an infimum. Use only the two
constructed values
\[
 L_A=\inf_{w\in A}\psi^A(\bar a,w),\qquad
 L_M=\inf_{b\in M}\psi^M(\bar a,b).
\]
For every $w\in A$,
$L_M\leq\psi^M(\bar a,w)\leq\psi^A(\bar a,w)+e$.
Therefore $L_M-e$ is a lower bound for all the values defining
$L_A$, and $L_M\leq L_A+e$.
For each $b\in M$, the witness contract supplies $w\in A$ with
\[
 L_A\leq\psi^A(\bar a,w)
      \leq\psi^M(\bar a,w)+e
      \leq\psi^M(\bar a,b)+\tau+e.
\]
Consequently $L_A-\tau-e$ is a lower bound for all the values
defining $L_M$, so $L_A\leq L_M+\tau+e$.
These two inequalities give the asserted error $e+\tau$.
No auxiliary infimum of $\psi^M$ over $A$ is assumed to exist.
The supremum argument is dual.
Induction gives \eqref{eq:cm-rank-budget}.
When the witness contract is available at every precision, let the
explicit tolerance tend to zero in the represented real values.
\end{proof}

This is weaker input than a metric net if a particular formula has an
efficient approximate minimizer even when the entire domain is not
totally bounded. It is also a more precise statement of the unresolved
obligation: one must produce and verify the lists in
\eqref{eq:cm-witness-test}. The theorem does not produce elementary
submodels from arbitrary existence claims, and its conclusion does
not supply a general first-order satisfaction decision procedure.

\begin{proposition}[Observable quotients]
\label{prop:cm-observable-quotient}
Suppose each sort has a projection $\pi_s:M_s\to K_s$ with a
specified section $j_s:K_s\to M_s$, so $\pi_sj_s$ is the identity.
Suppose every atomic recipe factors exactly through $\pi$:
\[
 P^M(\bar a)=P^K(\pi\bar a).
\]
Then every formula in the chosen language has
$\varphi^M(\bar a)=\varphi^K(\pi\bar a)$ whenever the values are
constructed. If $K$ has the compact presentation above, its finite
evaluator supplies these values even when $M$ itself is not compact.
\end{proposition}
\begin{proof}
For an atomic formula this is the factorization hypothesis.
Connectives preserve equality. For a quantifier the two sets of
subformula values coincide: projection sends every $M$-argument to a
$K$-argument, and the supplied section lifts every $K$-argument back.
Therefore their constructed suprema and infima coincide. Equivalently,
the graph of $\pi$ is a total correspondence of atomic error zero.
This proves the result without choosing any representative not already
provided by the section.
\end{proof}

For a constructive noncompact example, let $M$ be the closed unit ball
of $c_0$: its elements are real sequences $(x_j)$ with $|x_j|\leq1$,
supplied with a modulus $b_x(k)$ such that
$j\geq b_x(k)$ implies $|x_j|\leq2^{-k}$.
The supremum distance is a represented real, not an assumed supremum
of an arbitrary bounded set. To approximate it to precision $k$,
take the finite prefix ending before
$B=\max(1,b_x(k+2),b_y(k+2))$. At each of its coordinates,
approximate $x_j,y_j$ separately to error $2^{-(k+3)}$ and take
the absolute difference of the two rational samples. These are
nonnegative approximants to $|x_j-y_j|$ with error at most
$2^{-(k+2)}$. Their finite maximum
approximates the full distance to error at most $2^{-(k+1)}$:
outside the prefix every difference is at most $2^{-(k+1)}$,
and inside it the specified sampling bound applies. Rescaling by
$1/2$ gives diameter at most one if required.

Let $K=[-1,1]^d$. Projection takes the first $d$ coordinates; the
section pads with zeros and supplies the constant tail bound $d$.
The domain $M$ is not totally bounded: the coordinate unit vectors
are pairwise at supremum distance one, so a ball of radius less than
$1/2$ contains at most one of them. Nevertheless every formula whose
atoms depend only on the first $d$ coordinates reduces exactly to the
compact cube $K$. Rational coordinate formulas therefore have the
finite arithmetic interpretation above. Adding the full
supremum-distance predicate would invalidate the factorization
hypothesis. The scope of the observable language is part of the theorem.

\subsection{Exact formalization boundary and the remaining Lean work}
\label{subsec:cm-lean}

The supplied file \texttt{ContinuousSyntax.lean} defines arity-indexed
syntax \path{QuantitativeIPRA.Continuous.Formula}, a finite numerical
\texttt{Model}, environment extension, a finite evaluator, and
\texttt{Correspondence}. Its theorem
\path{QuantitativeIPRA.Continuous.formula_transfer} proves the
unchanged-error result for the full core basis, including arbitrary
alternation of finite infima and suprema. It uses the checked
\texttt{min\_correspondence} and \texttt{max\_correspondence}
lemmas in \texttt{FiniteExtrema.lean}. All carriers are nonempty ranges
$0,\ldots,\mathit{last}$. A rational score is represented by a natural
numerator over a common positive denominator. The predicate
\[
 \operatorname{Close}(e,a,b)\quad\equiv\quad
       a\leq b+e\ \wedge\ b\leq a+e
\]
avoids an incorrect use of truncated natural subtraction for absolute
error. Truncated complement agrees with $B-a$ on $0\leq a\leq B$.

These are checked Lean propositions about finite semantic structures,
\emph{not} terms of the repository's
$\texttt{Prf PRAAx}$ object theory. The full compact-presentation
compiler and its semantic adequacy theorem have the human proofs above
and the following explicit formalization path:
\begin{enumerate}
\item Introduce a record \texttt{CompactPresentation} containing the
      dense-code validity predicate, rational pseudometric, net lengths
      and entries, and object-theory derivations of
      \eqref{eq:cm-net}. Introduce \texttt{AtomicRecipe} with fixed
      $\PR$ terms for its evaluator and modulus, together with
      derivations of \eqref{eq:cm-atom-contract}. No field is a
      universal evaluator for arbitrary programs.
\item Compile each fixed formula to numerical $\PR$ terms using
      composition and the bounded folds already used for
      \texttt{ApproximateWitness}. For the core basis use one positive
      denominator, normalize rational constants, and compile complements
      as bounded subtraction.
\item Prove object-theory finite correspondence transport. The existing
      \texttt{FiniteNet}, \texttt{NetTransfer}, and
      \texttt{ApproximateWitness} modules supply net elimination and
      extremal-index infrastructure. Reproduce the four rational
      inequalities in the proof of
      Theorem~\ref{thm:cm-transfer}, then assemble one derivation per
      fixed syntax tree by the usual substitution and cut operations.
\item Define the finite relation
      $d_s(a,b)\leq2^{-U(h)}$ and prove its two coverage directions
      from \eqref{eq:cm-net}. Instantiate the object correspondence
      theorem to derive \eqref{eq:cm-cross-mesh}, and then
      \eqref{eq:cm-cauchy}. These are arithmetic endpoints with every
      mesh, precision, parameter and denominator explicit.
\item In a separate semantic namespace, define the real value by the
      regular Cauchy name. Prove \eqref{eq:cm-semantic-bounds} and
      Theorem~\ref{thm:cm-adequacy}, using the completion-extension
      theorem. A Mathlib realization may use completed metric spaces
      here, but that realization is an adequacy theorem and must not
      be relabeled an iPRA derivation.
\end{enumerate}

The fixed-formula proof transformation is enough for the proposed
applications. A uniform compiler on numerical syntax codes would
require its own primitive-recursive size bounds, proof-code encoding,
and correctness theorem. The checked Lean induction above is useful
evidence for that future implementation; it does not by itself supply
those additional object-theory results.

Finally, zero-valued truth remains different from finite observation.
The algorithm returns intervals of decreasing width containing
$\varphi^\CM(\bar a)$. It can certify a strict inequality when the
interval separates from its threshold. It need not decide the
statement $\varphi^\CM(\bar a)=0$. Neither equality of arbitrary real
names nor the complete theory of an arbitrary presented structure
has become decidable.

\section{A complex-analytic interface: power series and exact descent}\label{sec:complex}

Continuous gluing is not automatically holomorphic gluing. The weights in a partition of unity are usually not holomorphic, and averaging local holomorphic functions that disagree can destroy complex differentiability. We now state the extra data under which the continuous interface does preserve analytic structure. The construction uses finite Taylor sums, explicit tail bounds, and exact overlap equality expressed at every precision.

\subsection{A power-series presentation on an inner disk}

Identify \(\Q(i)\) with rational pairs and use the submultiplicative norm
\[
 |x+iy|_1=|x|+|y|.
\]
All calculations of this norm on rational pairs are PR. Its open balls are diamonds, which are also open in the usual complex topology. Let \(c\in\Q(i)\), \(R>0\), \(M>0\) be rational, and let \(a_n\in\Q(i)\) be a fixed PR coefficient sequence with a proved bound
\begin{equation}
 |a_n|_1R^n\leq M\qquad(n\in\N).
 \label{eq:complex-coeff}
\end{equation}
Fix a rational inner radius \(0<\rho<R\) and put \(\theta=\rho/R<1\). The finite evaluator is
\[
 P_N(z)=\sum_{n=0}^N a_n(z-c)^n.
\]
At rational \(z\) this is exact rational arithmetic. For \(|z-c|_1\leq\rho\), a finite geometric-sum calculation gives, for all \(m>N\),
\begin{equation}
 |P_m(z)-P_N(z)|_1
 \leq M\sum_{n=N+1}^{m}\theta^n
 \leq \frac{M\theta^{N+1}}{1-\theta}.
 \label{eq:complex-tail}
\end{equation}
This constructs a uniform Cauchy name without first assuming convergence of an infinite series.

The cutoff is explicitly PR. Choose natural \(H\geq1\) and \(b\geq0\) with
\[
 H(1-\theta)\geq1,\qquad M/(1-\theta)\leq2^b.
\]
The finite Bernoulli argument used for contractions gives \(\theta^H\leq1/2\). Thus
\[
  N(k)=H(k+b+1)
\]
makes the tail at most \(2^{-(k+1)}\). Both \(H\) and \(b\) can be chosen by bounded searches on the rational codes, or supplied as certificates. This is a PR sequence of finite polynomials representing a function \(f\) on the closed inner diamond.

\subsection{A derivative name and an explicit Taylor remainder}

The formal derivative polynomials, with \(N\geq1\), are
\(P_N'(z)=\sum_{n=1}^N n a_n(z-c)^{n-1}\). Put
\(\sigma=(1+\theta)/2\) and \(u=\theta/\sigma<1\). The finite estimate
\(n u^{n-1}\leq(1-u)^{-2}\) follows by bounding a term in the square of a finite geometric sum. Consequently
\begin{equation}
 \sum_{n>N}\frac{M}{R}n\theta^{n-1}
 \leq B_1\sigma^N,
 \qquad B_1=\frac{M}{R(1-u)^2(1-\sigma)}.
 \label{eq:derivative-tail}
\end{equation}
The notation on the left abbreviates the bound for every finite tail and its constructed limit. Choosing \(H_1(1-\sigma)\geq1\), \(B_1\leq2^{b_1}\), and \(N_1(k)=H_1(k+b_1+1)\) gives a PR Cauchy name \(g\) for the derivative series.

\begin{theorem}[Certified complex differentiability]\label{thm:complex-derivative}
For \(|z-c|_1+|h|_1\leq\rho\),
\begin{equation}
 |f(z+h)-f(z)-g(z)h|_1
 \leq T|h|_1^2,
 \qquad T=\frac{M}{R^2(1-\theta)^3}.
 \label{eq:complex-remainder}
\end{equation}
Thus \(f\) is complex differentiable at every point of the open inner diamond, with derivative \(g\). Given a positive margin to its boundary, the derivative condition has an explicit PR modulus. Also
\begin{equation}
 |f(z)|_1\leq\frac{M}{1-\theta},\qquad
 |f(z)-f(w)|_1\leq\frac{M}{R(1-\theta)^2}|z-w|_1
 \label{eq:complex-lip}
\end{equation}
for \(|z-c|_1,|w-c|_1\leq\rho\).
\end{theorem}

\begin{proof}
For \(n\geq2\), the binomial identity gives
\[
 |(v+h)^n-v^n-nv^{n-1}h|_1
 \leq |h|_1^2\sum_{j=2}^n \binom{n}{j}|v|_1^{n-j}|h|_1^{j-2}.
\]
Since
\(\binom{n}{j}=\frac{n(n-1)}{j(j-1)}\binom{n-2}{j-2}\leq\binom{n}{2}\binom{n-2}{j-2}\),
the sum is at most
\(\binom{n}{2}(|v|_1+|h|_1)^{n-2}\leq\binom{n}{2}\rho^{n-2}\).
Multiply by the coefficient bounds and sum over \(n\). Every finite partial sum is bounded by
\[
 \frac{M}{R^2}|h|_1^2
       \sum_{n=2}^{N}\binom{n}{2}\theta^{n-2}
 \leq\frac{M}{R^2(1-\theta)^3}|h|_1^2.
\]
The last inequality uses only finite sums: the coefficient of \(\theta^r\) for \(r\leq N\) in \((\sum_{j=0}^N\theta^j)^3\) is \(\binom{r+2}{2}\), and all omitted terms are nonnegative. The finite geometric sum is at most \((1-\theta)^{-1}\). The Cauchy bounds \eqref{eq:complex-tail} and \eqref{eq:derivative-tail} let the finite polynomial inequality pass to the represented limits.

For a requested derivative tolerance \(2^{-k}\), take \(|h|_1\leq2^{-k}/\max(1,T)\), together with the supplied inner-domain margin. Equation~\eqref{eq:complex-remainder} then bounds the linearization error by \(2^{-k}|h|_1\). This is the complex-linear derivative condition; no division by an undecidably zero increment is required.

There is also an explicit stability bound for the derivative approximations. Factoring the differences of powers in \(P_N'\) gives
\[
 |P_N'(z)-P_N'(w)|_1
 \leq\frac{M}{R^2}\sum_{n=2}^N n(n-1)\theta^{n-2}|z-w|_1
 \leq2T|z-w|_1.
\]
Together with \eqref{eq:derivative-tail}, this proves that the derivative name itself extends continuously to represented inputs, with Lipschitz constant \(2T\). The value name has the stability estimate proved next. Apply the completion extension to both names before transporting \eqref{eq:complex-remainder} from rational inputs to represented ones. Thus the notation \(g(z)\) at a represented point carries an actual evaluation procedure and modulus.

The value bound is the geometric estimate including the constant term. For the Lipschitz bound, factor
\(v^n-w^n=(v-w)\sum_{j=0}^{n-1}v^{n-1-j}w^j\).
Its norm is at most \(n\rho^{n-1}|v-w|_1\). The sum of the coefficient majorants is at most \(M/[R(1-\theta)^2]\), by the square-geometric-sum estimate. Passing through the uniform tails proves \eqref{eq:complex-lip}.
\end{proof}

The coefficient-majorant premise is substantial. A formal power series with no proved growth or tail bound is not automatically a continuous function. On the other hand, the premise is compatible with singularities outside the chosen disk and with arbitrary requested precision inside it. Rational coefficients already include many familiar analytic functions; coefficients given by regular names can also be used by allocating their finite-sum errors. At stage \(N\), bound the finite sum of absolute weights \(\sum_{n\leq N}|z-c|_1^n\), then query every coefficient deeply enough that this bound times the individual sample error fits half the output budget. The finite number of queries and their common precision are PR. Their coefficient coherence and majorant proofs remain explicit inputs.

\subsection{Exact analytic descent across finitely many charts}

Take \(N\geq1\) such series centered at \(c_i\), with inner radii \(r_i<R_i\). Let
\[
 w_i(z)=(r_i-|z-c_i|_1)_+,\qquad S(z)=\sum_iw_i(z).
\]
Choose \(\delta>0\), and use the certified subdomain \(D_0\) of rational points with \(S(z)\geq\delta\). Let \(\Omega=\{z:S(z)>\delta\}\). Strict membership in \(\Omega\) means a supplied positive rational margin, as usual for represented open sets. Every point of \(\Omega\) has rational names in \(D_0\); sufficiently accurate rational approximations remain there by the Lipschitz bound for \(S\).

Assume the local analytic functions agree on overlaps. The finite formulation is: at every rational \(z\) with positive weights in charts \(i,j\), the local sample values satisfy
\begin{equation}
 |f_i(p,z)-f_j(q,z)|_1\leq2^{-p}+2^{-q}.
 \label{eq:analytic-overlap}
\end{equation}
Uniform continuity extends this equality to represented overlap points. Let \(M_0\) and \(L_0\) be common rational bounds obtained from \eqref{eq:complex-lip}, with the sample-error allowance included in \(M_0\). The weighted gluing theorem applies.

\begin{corollary}[Holomorphic descent with quantitative guards]\label{cor:analytic-descent}
The weighted evaluation procedure defines a single represented function on the completion of \(D_0\). On \(\Omega\) it agrees with each active local analytic function and is holomorphic. A local derivative name and its remainder modulus are obtained from any chart with supplied positive clearance, using Theorem~\ref{thm:complex-derivative}. No chart is chosen by an exact sign test on a real number.
\end{corollary}

\begin{proof}
Equation~\eqref{eq:analytic-overlap} and the overlap estimate for the weighted sum show, at every requested precision, that the glued value agrees with each active chart. At \(z\in\Omega\), let \(\gamma>0\) witness \(S(z)\geq\delta+\gamma\). Since \(S\) is \(N\)-Lipschitz, increments of norm at most \(\gamma/(2N)\) stay inside the certified domain. At least one chart has positive weight with a rational lower bound; a finite evaluation of the weights with sufficient error slack supplies an index with weight at least \(\delta/(2N)\). For example, approximating each weight to error at most \(\delta/(8N)\) and scanning its rational lower endpoint finds such an index, since the maximum true weight is at least \(\delta/N\).

Increments smaller than \(\delta/(4N)\) stay in that chart. On this neighborhood the global function equals the chart function exactly. Its local derivative and remainder bound therefore apply to the global function. Combine the two domain margins with the derivative-error margin in \eqref{eq:complex-remainder}. These are explicit rational minima, so the resulting local modulus is PR. This proves holomorphicity on \(\Omega\), without differentiating the generally nonholomorphic weights.
\end{proof}

If overlap agreement is only within a fixed positive tolerance, the conclusion is only continuous approximate gluing. It does not imply holomorphicity. Nor does the theorem derive exact agreement merely from finitely many equal samples. In applications, agreement must come from a proved common expression, a controlled analytic identity, or another explicit all-precision argument. Branches with monodromy cannot be glued globally when this compatibility premise fails.

\subsection{Object formulas and the Lean route}

The new object-level endpoints are finite: the geometric tail inequality for two cutoffs; the derivative tail inequality; the polynomial Taylor remainder with a positive radius gap; and the weighted overlap estimate. Coefficients and sample points are coded rational pairs. The invariant for a partial polynomial tracks its value, derivative value, accumulated error, and geometric majorant. The three finite convolution sums used above are PR loops with cleared polynomial inequalities as quantifier-free invariants. The repository's rational and Gaussian-rational arithmetic modules provide the intended numerical backend.

In the semantic Lean layer, define \code{PowerSeriesPresentation} by its coefficient evaluator, coefficient bound, rational radii, and PR cutoff functions. Define \code{valueName}, \code{derivativeName}, and the exact proposition \code{TaylorRemainderGoal} using \eqref{eq:complex-remainder}. Prove the finite polynomial identity first; then the geometric majorants; then the Cauchy names; finally the derivative and descent interpretation. This ordering keeps the difficult analytic conclusion attached to explicit finite proof obligations. The blueprint supplied with the package states these definitions and targets; their completed \code{Prf PRAAx} proofs remain work to be implemented.

\section{An end-to-end problem over an infinite function space}\label{sec:variational}

To demonstrate that the framework is not restricted to finite combinatorics disguised by continuous terminology, consider the variational problem
\begin{equation}
  \sup\left\{\int_0^1 f(x)\,dx:
  f(0)=f(1)=0,\quad |f(x)-f(y)|\leq|x-y|\right\}.
 \label{eq:variational}
\end{equation}
The domain is an infinite-dimensional compact family of real-valued functions. The objective is integration. We will describe its presentation, obtain its exact value and maximizer, and display the finite arithmetic identity that proves the result.

\subsection{A compact presentation of the boundary condition}

Let \(t(x)=\min(x,1-x)\) on \([0,1]\). For any 1-Lipschitz \(f\) with the stated endpoint values, the two endpoint inequalities give
\begin{equation}
  -t(x)\leq f(x)\leq t(x).
 \label{eq:tent-bound}
\end{equation}
Conversely a 1-Lipschitz function obeying this bound has zero endpoints. Define
\[
  (Rf)(x)=\max\bigl(-t(x),\min(t(x),f(x))\bigr).
\]
The minimum and maximum of real-valued 1-Lipschitz functions are 1-Lipschitz. Hence \(R\) maps \(\mathcal F_{1,1}\) into the anchored class \(\mathcal F_0\) in \eqref{eq:variational}. It fixes every member of \(\mathcal F_0\). At each \(x\), clipping into the same interval is nonexpansive, so
\[
 D(Rf,Rg)\leq D(f,g).
\]
Apply \(R\) to every envelope center from Theorem~\ref{thm:function-net}. All resulting centers have the required exact boundary values and Lipschitz bound, and they form a net in \(\mathcal F_0\) with the same error. This avoids an unjustified claim that an arbitrary closed subset of an effectively compact space automatically has computable finite nets. Here the explicit nonexpansive retraction supplies the missing construction.

\subsection{Integration as a presented continuous functional}

For \(N\geq1\), let
\[
 S_N(f)=\frac1N\sum_{j=0}^{N-1}f(j/N).
\]
On each grid cell, the Lipschitz variation is at most \(1/N\). Thus any tagged refinement differs from the left-endpoint sum by at most \(1/N\). Comparing two sums on a common refinement shows
\begin{equation}
 |S_N(f)-S_{N'}(f)|\leq\frac1N+\frac1{N'}.
 \label{eq:riemann-name}
\end{equation}
Taking dyadic \(N\) produces a Cauchy name. Define \(I(f)\) to be this integral value. In usual Riemann semantics it is precisely \(\int_0^1f\). If every sampled value has rational error at most \(\eta\), the computed sum has an additional error at most \(\eta\), since the weights are nonnegative and sum to one. Therefore
\begin{equation}
 |\widetilde S_N(f)-I(f)|\leq N^{-1}+\eta.
 \label{eq:integral-error}
\end{equation}
Also \(|S_N(f)-S_N(g)|\leq D(f,g)\), and passage through \eqref{eq:riemann-name} gives
\begin{equation}
 |I(f)-I(g)|\leq D(f,g).
 \label{eq:integral-lip}
\end{equation}
This is a nonexpansive functional on the compact function sort. Its extrema can therefore be approximated by a finite maximum over the retracted envelope net. The error is the net radius plus the integration-evaluation error. No choice of an exact optimizing function is needed to compute the optimal value.

All the arithmetic content of \eqref{eq:riemann-name} is finite: a common rational refinement, a bound for each cell, and the sum of the cell lengths. Prove the inequality for the first \(j\) cells, clearing the common denominator, and induct on \(j\). This is the quantifier-free sum invariant already exemplified by the repository's grid and refinement modules. The integral symbol itself never appears in the object-language proof.

\subsection{Exact optimum and an explicit attainer}

\begin{theorem}[The tent variational problem]\label{thm:tent}
The value of \eqref{eq:variational} is \(1/4\). It is attained by
\(t(x)=\min(x,1-x)\). The finite rational inequalities proving the upper bound and the dyadic identity proving its value are derivable in \ipra\ using quantifier-free induction. The semantic integral statement follows from the preceding integral-name construction.
\end{theorem}

\begin{proof}
The function \(t\) is 1-Lipschitz, is nonnegative, and has zero endpoint values. Equation~\eqref{eq:tent-bound} gives \(f(x)\leq t(x)\) for every admissible \(f\). For \(N=2n\), \(n\geq1\), a direct finite computation is
\begin{align*}
 \sum_{j=0}^{2n-1}\min(j,2n-j)
 &=\sum_{j=0}^{n}j+\sum_{j=1}^{n-1}j\\
 &=\frac{n(n+1)}2+\frac{n(n-1)}2=n^2.
\end{align*}
Thus \(S_{2n}(t)=n^2/(2n)^2=1/4\). Summing \(f(j/N)\leq t(j/N)\) gives \(S_{2n}(f)\leq1/4\). With rational samples of error \(\eta\), the finite conclusion is instead
\begin{equation}
   \widetilde S_{2n}(f)\leq\frac14+\eta.
 \label{eq:finite-variational}
\end{equation}
The integral estimate gives \(I(f)\leq1/4+1/(2n)+2\eta\). At arbitrary requested precision choose explicit dyadic \(n\) and \(\eta\) making the last two terms as small as desired. This is exactly the inequality \(I(f)\leq1/4\) between represented reals. The constant dyadic sums for \(t\) show \(I(t)=1/4\), proving attainment.

For the arithmetic derivation encode the sample table \(a_j\), its denominator data, and an error numerator. The finite premises are the endpoint-distance bounds that imply \(a_j\leq\min(j/N,1-j/N)+\eta\). Multiply by a positive common denominator. The invariant for the first \(r\) entries is the partial-sum inequality plus \(r\eta\); it is a numerical quantifier-free inequality. The displayed sum-of-integers identities follow from quantifier-free induction. This yields \eqref{eq:finite-variational} as a first-order theorem for arbitrary finite tables satisfying the specified contract. The presentation provider proves those premises for any named admissible \(f\).
\end{proof}

\subsection{A quantitative uniqueness estimate}

The exact optimizer is also recoverable by the residual-diameter method, even if the search algorithm never recognizes the tent function. Let
\[
 J(f)=\frac14-I(f)=\int_0^1(t-f),\qquad f\in\mathcal F_0.
\]
The difference \(g=t-f\) is nonnegative and 2-Lipschitz. At a point where \(g(x)=a\), Lipschitz continuity gives
\[
 g(y)\geq\max(0,a-2|x-y|).
\]
Since \(g(0)=g(1)=0\), we have \(a\leq2\min(x,1-x)\). Consequently the entire triangular interval \([x-a/2,x+a/2]\) lies in \([0,1]\). Its area is \(a^2/2\). Integrating the lower bound, or using explicit dyadic lower sums and their error bound, gives
\begin{equation}
   g(x)^2\leq2J(f),\qquad
   D(f,t)^2\leq2J(f).
 \label{eq:tent-residual}
\end{equation}
The supremum in the second inequality needs no attained point: apply the first inequality to every point or to finite net maxima, then pass through their error bounds.

If \(J(f),J(h)\leq2^{-(2k+3)}\), each of \(D(f,t),D(h,t)\) is at most \(2^{-(k+1)}\), and hence \(D(f,h)\leq2^{-k}\). Thus a residual-diameter modulus is
\[
  \kappa(k)=2k+3.
\]
The square-root notation is unnecessary in the finite arithmetic: comparisons are made after squaring nonnegative rational quantities. An arbitrary PR sequence of approximate maximizers with this residual schedule therefore converges coherently to the tent function. This connects continuous optimization, quantitative uniqueness, compact function sorts, and the completion interface in one example.

For completeness, the integral of the lower triangular hat is itself a finite rational certificate. If \(x,a\) are rational and \(0\leq a\leq2\min(x,1-x)\), partition the interval at \(0,x-a/2,x,x+a/2,1\). The hat is affine on every piece. Its two nonzero trapezoid areas are each \((a/2)(0+a)/2=a^2/4\), so their sum is \(a^2/2\). The trapezoid identity for an affine function follows by summing an arithmetic progression; it is a quantifier-free arithmetic theorem after denominators are cleared. Since the hat is 2-Lipschitz, any uniform-grid sum differs from this exact piecewise-affine value by at most \(2/N\), by the common-refinement proof of \eqref{eq:riemann-name}. Summing the pointwise lower bounds and sending this explicit error to zero proves \eqref{eq:tent-residual} for rational centers and heights.

For a represented center \(x\) and height \(a=g(x)\), choose a rational center \(r\in[0,1]\) with \(|r-x|\leq\epsilon\) (clipping a rational approximation to the interval preserves its error) and a rational \(b\) with \(|b-a|\leq\epsilon\). The nonnegative rational height \(c=\max(0,b-3\epsilon)\) satisfies \(c\leq g(r)\), because \(g\) is 2-Lipschitz. It also satisfies \(c\leq2\min(r,1-r)\) by the endpoint bounds. Apply the rational-center estimate to the hat of height \(c\), and let the explicit approximation error tend to zero. This proves the assertion for every represented point without integration over an unspecified real partition. The estimate is sharp: for \(f=-t\), \(D(f,t)=1\) and \(J(f)=1/2\), giving equality. Neither exact maximizer selection nor compactness choice is needed.

\section{Lean artifacts and the exact proof boundary}\label{sec:verification}

\subsection{What was compiled}

The package contains a standalone Lean project under \path{lean/}. It was built in this session with Lean 4.32.0, compiler commit
\path{8c9756b28d64dab099da31a4c09229a9e6a2ef35}. It uses only Lean's bundled standard library; it does not require Mathlib or a ProveIt checkout. Its five modules contain 41 named proved theorems. The default Lake build has seven jobs, all successful.

\begin{longtable}{@{}P{.27\textwidth}P{.65\textwidth}@{}}
\toprule File & Checked content\\\midrule\endhead
\path{FiniteExtrema.lean} & Actual nonempty min/max folds; attained argmin/argmax; least-minimizer property; two-sided error estimates; total-correspondence transport.\\
\path{CompactQuantifier.lean} & Finite-net lower/upper bounds, selected approximate witnesses, and comparison between two nets with different mesh and sample errors.\\
\path{ContinuousSyntax.lean} & Single-sort, arity-indexed formula syntax, finite model semantics, binder environments, total correspondence, and formula transfer for the complete stated connective/quantifier basis.\\
\path{Examples.lean} & An endpoint/midpoint tent-function model comparison that attains the transfer error bound.\\
\path{Audit.lean} & Axiom reports for the main exported theorems.\\
\bottomrule
\end{longtable}

The principal declaration is
\begin{center}
\path{QuantitativeIPRA.Continuous.formula_transfer}.
\end{center}
The checked theorem is the single-sort finite semantic core; its two model carrier sizes may differ. It assumes a two-sided total relation and atomic closeness for all related tuples. Its conclusion is closeness by the same natural score error for the evaluation of every arity-correct formula. The binders genuinely extend the variable environment. This is not a placeholder that assumes formula-level correctness in a record field. The many-sorted version proved in the model-theory chapter remains a Lean generalization: replace the one carrier by a family indexed by sorts, index variable environments and atomic argument lists by their sort lists, and give each binder its sort. The same induction then uses totality for the sort bound by that binder. Formulas mixing function and point sorts require this generalization.

The arithmetic error predicate is
\[
 \mathsf{Close}(e,a,b)\quad\equiv\quad a\leq b+e\ \wedge\ b\leq a+e.
\]
For a known positive common denominator \(D\), it states an absolute rational error of at most \(e/D\). Natural truncated complement is nonexpansive on all natural scores, and agrees with the intended real complement after scaling on the prescribed interval \(0\leq a\leq D\). Every finite carrier is indexed by \(0,\ldots,n\), avoiding an undefined empty extremum.

The reported dependencies of the main finite transfer theorem are \code{propext} and \code{Quot.sound}. The audit has no \code{Classical.choice}, \code{sorryAx}, or custom correctness axiom. These facts concern the exported Lean theorem dependencies. They do not convert the semantic Lean theorem into an object-language \ipra\ proof.

\subsection{What remains an outline}

There are no new completed \code{Prf PRAAx} proof values for the continuous framework in this package. The complete human proofs specify the numerical inequalities and primitive-recursive schedules, but their deep reification is still to be implemented. This boundary is made explicit in the source layout:

\begin{itemize}
\item \path{lean/} contains the compiled standalone semantic kernel and its logs.
\item \path{lean/ObjectTheoryPlan.lean.txt} is a detailed plan for the ordinary ProveIt calculus, with inspected import paths, exact existing API names, intended object terms, formulas, and proof construction.
\item \path{blueprints/AnalyticTopologyBlueprint.lean} defines semantic records and propositions named \code{...Goal}; it is an uncompiled blueprint, and those propositions are not asserted theorems.
\item \path{blueprints/FunctionSpaceAnalyticPlan.lean.txt} supplies the remaining function-space and complex-analysis definitions, exact targets, and dependency order. It is also uncompiled.
\end{itemize}

\subsection{The shortest route to actual ordinary-iPRA endpoints}

The initial object-theory implementation should add two finite correspondence lemmas, rather than start by formalizing all of real analysis. Fix PR score programs \(f,g\), a relation bit \(R\), and numerical parameters \(x\). Let \(N,M\) be the last carrier indices and \(e\) a natural error. The input formulas are
\begin{align*}
 &\forall i\leq N\;\exists j\leq M\quad R(i,j,x)=1,\\
 &\forall j\leq M\;\exists i\leq N\quad R(i,j,x)=1,\\
 &\forall i\leq N\;\forall j\leq M\quad
 R(i,j,x)=1\ \longrightarrow\ \mathsf{Close}(e,f(i,x),g(j,x)).
\end{align*}
The intended conclusion is \(\mathsf{Close}(e,\min_i f(i,x),\min_jg(j,x))\), and similarly for maxima. These are fully first-order formulas. The inputs can be compiled through the existing \path{BoundedSyntax.BFm} interface.

The proof uses the actual least argmin. Instantiate one totality premise at the other carrier's minimizing index, eliminate its bounded witness, then combine the pointwise error and minimum specification. Repeat in the other direction. The \path{ApproximateWitness}, \path{RationalArgmin}, \path{FiniteQueryDemand}, and \path{BoundedSyntaxInduction} APIs already provide the needed scan and bounded-induction infrastructure. The supplied plan names their actual declarations at the pinned revision.

Next, perform external structural recursion on one fixed formula. Generate both a PR evaluator term and an ordinary proof tree. Min, max, and complement use arithmetic lemmas; the quantifier steps use the new correspondence lemmas; substitution inserts the coded extended environment. Since this is recursion in the metatheory over fixed syntax, it does not invoke a universal numerical evaluator or induction on a truth predicate.

Only then add the presentation adapters: actual net codes and coverage proofs, atomic evaluator and modulus proofs, a decidable relation with the named-distance slack, and the cross-mesh calculation. Their result is a literal all-precision Cauchy sentence about a fixed PR function. The semantic real-valued interpretation is proved in a separate namespace.

For the other chapters the dependency order is equally concrete:
\begin{enumerate}
\item Derive the rational small-error, finite reindexing, and completion-equivalence lemmas for fixed name programs. Instantiate all-precision hypotheses at explicit indices; do not use them as induction predicates.
\item Derive one weighted-summand inequality, then its finite-sum invariant and positive-denominator cancellation. Add safe inner weights for approximate metrics.
\item Connect the existing coherence endpoint to the residual's canonical evaluation. The proof must output equality of represented residual names, not only a bound on distances between approximate solutions.
\item Implement the rational envelope table evaluator, its Lipschitz inequality, the finite grid cover, and the table quantizer. Each selected center is proved legal directly. Then derive the supremum-distance-name bound.
\item Implement finite polynomial value/derivative sums and their geometric majorants. Derive the Taylor remainder before claiming complex differentiability or holomorphic descent.
\item Specialize these interfaces to the cubic inverse and the tent variational problem. Retain all domain, denominator, modulus, and approximate-solvability hypotheses until a specific provider discharges them.
\end{enumerate}

The formalization is complete only when these endpoints have actual \code{Prf PRAAx} values, compile in the pinned project or an explicitly migrated revision, and pass an axiom audit. A theorem about \code{PR.ev} in Lean's ambient mathematics is useful for adequacy but does not satisfy that gate.

\subsection{Executable evidence and reproducibility}

The exact-rational Python code covers envelope selection, supremum-distance observations, finite continuous formulas, the tent finite-sum certificate, power-series cutoff bounds, cube-root grid search, and weighted gluing. The continuous validation script performed 13,710 finite checks. The second example runner records known exact values, sharp examples, and computed error allowances. All use exact fractions rather than floating-point comparisons. These runs validate implementations and constants on the recorded examples; they are not universal proofs or formal proof certificates.

To reproduce the standalone kernel, enter \path{lean/} and run \code{lake build}; then run \code{lake env lean Audit.lean}. To reproduce the arithmetic examples, run \code{python3 code/run\_examples.py} and \code{python3 code/validate\_continuous.py} from the package root. To rebuild the PDF, run \code{./build\_pdf.sh}. The source manifest records the inspected repository files and content hashes. No GitHub branch or repository contents were modified during this research.

\section{Resolved difficulties, retained hypotheses, and limits}\label{sec:limits}

\subsection{What the framework resolves}

\begin{longtable}{@{}P{.25\textwidth}P{.35\textwidth}P{.32\textwidth}@{}}
\toprule Difficulty & Positive result & Data still required\\\midrule\endhead
Finite solutions at each scale need not form a point. & Residual diameter control gives a Cauchy name, canonical residual value zero, and uniqueness in the completion. & An actual PR approximate-solution routine and a proved PR diameter modulus.\\
A Cauchy estimate does not identify the value of a function at the limit. & Quantitative extension, density, and uniqueness identify the completed evaluation. & A proved modulus and extensional finite evaluator.\\
Local chart selection can be discontinuous or undecidable. & Weighted gluing is continuous with constant \(L+2NM/\delta\); safe inner weights handle named distances. & A positive denominator certificate and local correctness on every used support.\\
Finite sample tables may not be realizable as legitimate functions. & Every clipped envelope is a bounded Lipschitz function; quantized tables give an explicit net. & Shared range/continuity bounds; evaluator data when selecting a center near an input function.\\
Quantification ranges over continuous objects and functions. & Total correspondences yield finite continuous-formula values with no quantifier-depth error loss for the core connectives. & Compact presentations, atomic moduli, and finite evaluators.\\
Derivative separation fails at a singular inverse. & A H\"older residual modulus constructs the real cube root through zero. & Specific cubic separation and approximate-solvability estimates.\\
Continuous partitions of unity do not preserve holomorphicity. & Exact overlap equality makes the glued function locally equal to a certified analytic chart. & Power-series majorants, positive chart/domain margins, and all-precision overlap equality.\\
An infinite-domain optimization problem seems to demand choice. & The tent problem has exact value \(1/4\), an explicit attainer, and residual modulus \(2k+3\). & A concrete retraction and quantitative uniqueness proof; these are supplied for this example.\\
\bottomrule
\end{longtable}

\subsection{What is not implied}

The results do not turn every true theorem of analysis into an \ipra\ theorem. Qualitative continuity does not supply a PR modulus. A true existential statement does not supply a PR witness or a PR search bound. The logical complexity of a finite verification predicate must be addressed by actual bounded Boolean compilation, and the loop invariant must be proved with quantifier-free induction. These are mathematical obligations, not implementation details to omit.

Neither total boundedness nor finite solvability at every scale selects a computable infinite branch in general. Our function-space construction supplies finite nets and a completion operation for already modulated Cauchy sequences; it does not extract subsequences from arbitrary sequences. Computing a supremum value is distinct from computing a point attaining it. The residual theorem adds exactly the quantitative uniqueness information used to bridge that distinction in the stated examples.

The model-theoretic results evaluate a fixed formula on a supplied structure or compare supplied structures. They do not construct a model of an arbitrary consistent theory, prove first-order compactness inside \ipra, supply ultraproducts, or decide exact zero-valued truth. The observable-quotient example works only for a language whose atoms factor through the specified projection; adding the full supremum norm of the original noncompact space invalidates that premise.

The categorical results are equally scoped. Completion has a universal extension property in the external category of the chosen presentations. Finite correspondences compose with an additive error. Fixed Lipschitz bounds support the stated evaluation/currying correspondence. These facts do not assert a category of all sets or all structures inside first-order arithmetic, unrestricted exponentials, arbitrary quotients, or arbitrary colimits. A setoid interpretation needs derivations of the translated source axioms, as the existing ProveIt interface already requires.

Finally, the successful singular real cube root does not remove the known obstruction to uniform branch selection. The complex cube root has no single holomorphic inverse around zero, and monodromy can prevent global analytic gluing even when every chart is valid. Our analytic descent theorem retains the exact overlap premise rather than inferring it from a finite cover.

\subsection{A practical interface for future packets}

An analytic or abstract packet fits this framework when it explicitly supplies five components:
\begin{enumerate}
\item A numerical presentation of the relevant objects and finite observations, including domain and denominator guards.
\item Fixed PR evaluators or finite-table algorithms, with moduli, loop bounds, and nonempty finite supports where needed.
\item Finite correctness and stability inequalities, with quantifier-free or Booleanized bounded invariants.
\item A coherence or representation-independence argument, including any diameter, separation, or overlap modulus.
\item An adequacy theorem identifying the resulting names or formula values with the desired continuous meaning.
\end{enumerate}

This checklist does not add a new axiom. It locates the mathematical work and makes reuse possible. For the present report the most valuable next implementation is the ordinary-iPRA finite correspondence theorem, followed by the bridge from the existing coherent-minimizer name to its canonical completed residual. The compiled finite kernel and the explicit numerical schedules make those next steps concrete without claiming that the remaining deep formalization has already been completed.

\begin{thebibliography}{99}
\bibitem{repo} V. Reshetnikov and contributors, \emph{ProveIt: Primitive Recursive Arithmetic}, source baseline \path{5f3be02388442f16487d2dc5c2529961dd783d2b}, 20 November 2026. \url{https://github.com/VladimirReshetnikov/ProveIt/tree/5f3be02388442f16487d2dc5c2529961dd783d2b/Logic/PrimitiveRecursiveArithmetic}. Exact inspected files appear in \path{metadata/source_manifest.json}.
\bibitem{oconnor} R. O'Connor, \emph{A monadic, functional implementation of real numbers}, Mathematical Structures in Computer Science 17 (2007), 129--159. \url{https://doi.org/10.1017/S0960129506005871}. Author's text: \url{https://r6.ca/FewDigits/CReal20060514.pdf}.
\bibitem{continuouslogic} I. Ben Yaacov, A. Berenstein, C. W. Henson, A. Usvyatsov, \emph{Model theory for metric structures}, in Model Theory with Applications to Algebra and Analysis, vol.~2, London Mathematical Society Lecture Note Series 350 (2008), 315--427. Author's text: \url{https://math.univ-lyon1.fr/~begnac/articles/mtfms.pdf}.
\bibitem{kohlenbach} U. Kohlenbach, \emph{Things that can and things that can't be done in PRA}, Annals of Pure and Applied Logic 102 (2000), 223--245; BRICS Report RS-98-18 (1998). \url{https://www2.mathematik.tu-darmstadt.de/~kohlenbach/paper18.pdf}.
\bibitem{mcshane} E. J. McShane, \emph{Extension of range of functions}, Bulletin of the American Mathematical Society 40 (1934), 837--842. \url{https://doi.org/10.1090/S0002-9904-1934-05978-0}.
\bibitem{bishop} E. Bishop and D. Bridges, \emph{Constructive Analysis}, Grundlehren der mathematischen Wissenschaften 279, Springer, 1985. \url{https://doi.org/10.1007/978-3-642-61667-9}.
\end{thebibliography}
\end{document}
